\section{Arithmetic factors, operators, and positive realizations}
\label{syn-arith:section}

An Euler factor can arise from an integral or from an operator trace.
All Hilbert inner products below are conjugate-linear in the first variable.
For a prime $p$, the elementary equality is
\begin{equation}
 \int_{\mathbb Q_p^\times}\mathbf1_{\mathbb Z_p}(x)|x|_p^s\,d^\times x
 =\operatorname{Tr}_{\ell^2(\mathbb N_0)}e^{-sN_p}
 =\frac1{1-p^{-s}},\qquad \operatorname{Re}s>0.
 \label{syn-arith:spine-local}
\end{equation}
The measure has $\operatorname{vol}(\mathbb Z_p^\times)=1$ and
$N_pe_r=r\log p\,e_r$.  The integral and trace therefore agree after
their normalizations and domains have been fixed.
An arithmetic cohomology theory would require substantially more data.
In particular, the same energy space supports algebras with different bar homology.

There are two further exact bridges.  Chinese remainders turn finite
congruence conditions into tensor products and positive sieve operators.
Weil's explicit formula turns the Riemann hypothesis into positivity of
one specified convolution functional.  The sieve operators are positive
before any prime-distribution estimate.  The positivity of the Weil
functional is itself the unresolved arithmetic assertion.

\subsection{Local Mellin integrals and the completed zeta function}
\label{syn-arith:tate}

Write $\mathbb A=\mathbb A_{\mathbb Q}$ and $|p|_p=p^{-1}$.
Fix the additive character
\[
 \psi(x)=\exp\left(2\pi i\left(x_\infty-
                  \sum_p\{x_p\}_p\right)\right)
 \quad\text{on }\mathbb A/\mathbb Q.
\]
Here $\{x_p\}_p$ is the $p$-primary fractional part.
Self-dual additive measures satisfy $\operatorname{vol}(\mathbb Z_p)=1$
and give Lebesgue measure at infinity.  Set
\[
 d^\times x_p=(1-p^{-1})^{-1}\frac{dx_p}{|x_p|_p},\qquad
 d^\times x_\infty=\frac{dx}{|x|},\qquad
 \Gamma_{\mathbb R}(s)=\pi^{-s/2}\Gamma(s/2).
\]
For a Schwartz--Bruhat function $\phi_v$ on $\mathbb Q_v$, put
$Z_v(\phi_v,s)=\int\phi_v(x)|x|_v^s\,d^\times x$ wherever this is absolutely convergent.

\begin{proposition}[Local Tate factors]
\label{syn-arith:local-tate}
For $\operatorname{Re}s>0$,
\[
 Z_p(\mathbf1_{\mathbb Z_p},s)=(1-p^{-s})^{-1},\qquad
 Z_\infty(e^{-\pi x^2},s)=\Gamma_{\mathbb R}(s).
\]
With additive measure the finite integral is
\[
 \int_{\mathbb Z_p}|x|_p^{s-1}\,dx_p
 =\frac{1-p^{-1}}{1-p^{-s}}.
\]
\end{proposition}
\begin{proof}
The annulus $p^r\mathbb Z_p^\times$ has multiplicative volume one.
Only $r\geq0$ contributes, giving $\sum_{r\geq0}p^{-rs}$.
Its additive volume is $(1-p^{-1})p^{-r}$, which gives the second normalization.
At infinity, splitting into two half-lines and substituting $u=\pi x^2$ gives
\[
 2\int_0^\infty e^{-\pi x^2}x^{s-1}\,dx
 =\pi^{-s/2}\int_0^\infty e^{-u}u^{s/2-1}\,du.
\]
The endpoint at zero requires exactly $\operatorname{Re}s>0$.
\end{proof}

For $\phi=e^{-\pi x_\infty^2}\otimes\bigotimes_p\mathbf1_{\mathbb Z_p}$,
Tate's adelic integral is
\[
 Z(\phi,s)=\int_{\mathbb A^\times}\phi(x)|x|_{\mathbb A}^s\,d^\times x,
 \qquad |x|_{\mathbb A}=|x_\infty|\prod_p|x_p|_p.
\]
Absolute convergence holds for $\operatorname{Re}s>1$.  Restricted-product
integration and Proposition~\ref{syn-arith:local-tate} give
\[
 Z(\phi,s)=\Gamma_{\mathbb R}(s)\prod_p(1-p^{-s})^{-1}
          =\Gamma_{\mathbb R}(s)\zeta(s)=:\Lambda(s).
\]
These conventions agree with Tate's local and global calculation
\cite[\S\S2.2--2.5, pp.~316--321; \S4.4, Main Theorem~4.4.1, p.~339]{syn-arith:Tate}.

\begin{theorem}[Continuation from the theta integral]
\label{syn-arith:theta}
The function $\Lambda$ extends meromorphically to $\mathbb C$,
with residues $-1$ at $0$ and $1$ at $1$.  It satisfies
$\Lambda(s)=\Lambda(1-s)$.  Consequently
\[
 \xi(s)=\tfrac12s(s-1)\Lambda(s)
\]
is entire and satisfies $\xi(s)=\xi(1-s)$.
\end{theorem}
\begin{proof}
Let $\theta(t)=\sum_{n\in\mathbb Z}e^{-\pi n^2t}$ for $t>0$.
Poisson summation for the Gaussian gives
$\theta(t)=t^{-1/2}\theta(t^{-1})$.  Termwise integration in
$\operatorname{Re}s>1$ gives
\[
 \Lambda(s)=\frac12\int_0^\infty(\theta(t)-1)t^{s/2-1}\,dt.
\]
Split at $1$, use the theta identity on $(0,1)$, and substitute $t=u^{-1}$.
The result is
\begin{equation}
 \Lambda(s)=\frac1{s-1}-\frac1s+
 \frac12\int_1^\infty(\theta(t)-1)
       \left(t^{s/2-1}+t^{(1-s)/2-1}\right)dt.
 \label{syn-arith:theta-continuation}
\end{equation}
The integral is entire because $\theta(t)-1$ decays exponentially.
Both the rational terms and the integral are invariant under $s\mapsto1-s$.
The residues and the assertion about $\xi$ follow.
\end{proof}

For a primitive Dirichlet character $\chi$ of conductor $q$, let
$a\in\{0,1\}$ satisfy $\chi(-1)=(-1)^a$.  The corresponding scalar completion is
\[
 \Lambda(s,\chi)=(q/\pi)^{(s+a)/2}\Gamma((s+a)/2)L(s,\chi).
\]
Tate's functional equation gives
$\Lambda(s,\chi)=\varepsilon(\chi)\Lambda(1-s,\overline\chi)$,
where $\varepsilon(\chi)=i^{-a}q^{-1/2}\sum_{r\bmod q}\chi(r)e^{2\pi ir/q}$.
Replacing a single Euler factor leaves the conductor and infinity factor unspecified.

For a number field $K$ of signature $(r_1,r_2)$, discriminant $d_K$, class
number $h_K$, regulator $R_K$, and number $w_K$ of roots of unity, the
analytic class-number formula is
\[
 \operatorname*{Res}_{s=1}\zeta_K(s)
 =\frac{2^{r_1}(2\pi)^{r_2}h_KR_K}{w_K\sqrt{|d_K|}}.
\]
It is the general residue calculation of Tate's global zeta integral
\cite[\S4.4, pp.~338--339]{syn-arith:Tate}; the regulator of a rank-zero unit group
is one.  The Minkowski lattice $\mathcal O_K\subset\mathbb R^{r_1}\times
\mathbb C^{r_2}$ has covolume $2^{-r_2}\sqrt{|d_K|}$ for ordinary Lebesgue
measure on each real and complex coordinate, and $\sqrt{|d_K|}$ when each
complex measure is multiplied by two.  This follows by comparing the real
embedding matrix determinant with the square root of the discriminant.
Neither the logarithm of this covolume nor the zeta residue is a vector-space
dimension.  Over $\mathbb Q$ the residue is one; the number two obtained
by multiplying it by two is not a calculation of a conformal central charge.

\subsection{Prime equilibrium states and the global trace threshold}
\label{syn-arith:operators}

On $\mathscr H_p=\ell^2(\mathbb N_0)$, let $S_pe_r=e_{r+1}$ and
$\mathcal T_p=C^*(S_p)$.  Give $N_p$ its maximal domain
\[
 \mathcal D(N_p)=\left\{(a_r):\sum_{r\geq0}r^2(\log p)^2|a_r|^2<\infty\right\}.
\]
The diagonal operator is positive and self-adjoint.  Finite-support vectors form a core.
The time evolution is $\alpha_{p,t}(a)=e^{itN_p}ae^{-itN_p}$.
In particular, $\alpha_{p,t}(S_p)=p^{it}S_p$.
For inverse temperature $\beta>0$, the Kubo--Martin--Schwinger (KMS)
convention is $\varphi(ab)=\varphi(b\alpha_{i\beta}(a))$ on analytic elements.

\begin{proposition}[Prime Gibbs state]
\label{syn-arith:prime-kms}
The operator $e^{-sN_p}$ is trace class exactly for $\operatorname{Re}s>0$.
For every $\beta>0$, the unique KMS$_\beta$ state on $\mathcal T_p$ is
\[
 \varphi_{\beta,p}(a)=(1-p^{-\beta})\operatorname{Tr}(e^{-\beta N_p}a),
 \qquad
 \varphi_{\beta,p}(S_p^mS_p^{*n})=\delta_{mn}p^{-\beta m}.
\]
\end{proposition}
\begin{proof}
For $\operatorname{Re}s\geq0$, the singular values are $p^{-r\operatorname{Re}s}$.
They are summable exactly in the stated half-plane.
For $\operatorname{Re}s<0$ the diagonal exponential is unbounded.
Writing $q=p^{-\beta}$, the density $(1-q)e^{-\beta N_p}$ has trace one.
The monomials $S_p^mS_p^{*n}$ span a norm-dense algebra of entire analytic elements.
The identity
$e^{-\beta N_p}a=\alpha_{p,i\beta}(a)e^{-\beta N_p}$
and trace cyclicity prove KMS on this algebra, hence on $\mathcal T_p$.
Invariance annihilates every monomial with $m\ne n$.
For the others, KMS gives
\[
 \varphi(S_p^mS_p^{*m})=p^{-\beta m}\varphi(S_p^{*m}S_p^m)=p^{-\beta m}.
\]
This determines a state uniquely by density.
\end{proof}

For finite prime sets $F$, take minimal $C^*$ tensor products
$\mathcal T_F=\bigotimes_{p\in F}^{\min}\mathcal T_p$ with unit-insertion maps.
Their norm inductive limit is $\mathcal T_{\mathrm{fin}}$.
The Hilbert tensor product relative to $e_{p,0}$ identifies with
$\ell^2(\mathbb N^\times)$ by
\[
 \bigotimes_pe_{p,r_p}\longmapsto e_{\prod_pp^{r_p}}.
\]
Every elementary tensor has finite support in $p$.
On this space put
\[
 He_n=(\log n)e_n,\qquad
 \mathcal D(H)=\left\{(a_n):\sum_{n\geq1}(\log n)^2|a_n|^2<\infty\right\}.
\]

\begin{theorem}[Global prime tensor system]
\label{syn-arith:global-kms}
For every $\beta>0$, the product state
$\Phi_\beta=\bigotimes_p\varphi_{\beta,p}$ is the unique KMS$_\beta$ state
on $\mathcal T_{\mathrm{fin}}$.  On the vacuum tensor representation,
\begin{equation}
 e^{-sH}\text{ is trace class}\quad\Longleftrightarrow\quad
 \operatorname{Re}s>1,
 \qquad \operatorname{Tr}e^{-sH}=\zeta(s).
 \label{syn-arith:trace-threshold}
\end{equation}
For $\beta>1$, its Gibbs density is $\zeta(\beta)^{-1}e^{-\beta H}$.
For $0<\beta\leq1$, the product state exists but has no normal extension
to $\mathcal B(\ell^2(\mathbb N^\times))$.
\end{theorem}
\begin{proof}
The local product states are compatible under unit insertion.
On a finite tensor product, a monomial has frequency
$\sum_p(m_p-n_p)\log p$.  Unique factorization makes this zero only when
$m_p=n_p$ for every $p$.  KMS then forces its expectation to be
$\prod_pp^{-\beta m_p}$.  This proves uniqueness at finite stages and at the limit.
The singular-value sum for the global exponential is
$\sum_{n\geq1}n^{-\operatorname{Re}s}$, proving
\eqref{syn-arith:trace-threshold} and the Gibbs formula.

For the last assertion, let
$P_F=\prod_{p\in F}(1-S_pS_p^*)$ in the vacuum representation.
It projects onto integers divisible by none of the primes in $F$.
As $F$ increases, $P_F$ decreases strongly to the rank-one vacuum projection $P_1$.
Its expectation is $\prod_{p\in F}(1-p^{-\beta})$, which tends to zero
for $\beta\leq1$.  A normal extension would therefore give $\Phi_\beta(P_1)=0$.
For $S_n=\prod_pS_p^{v_p(n)}$, normality also gives
\[
 \Phi_\beta(S_nP_1S_n^*)
 =\lim_F n^{-\beta}\Phi_\beta(P_F)=0.
\]
These rank-one projections sum strongly to the identity.
Countable additivity would give $\Phi_\beta(1)=0$, a contradiction.
\end{proof}

The full Bost--Connes algebra additionally contains commuting unitaries
$e(r)$ for $r\in\mathbb Q/\mathbb Z$.  Its defining relations include
\[
\begin{gathered}
 \mu_m\mu_n=\mu_{mn},\quad \mu_n^*\mu_n=1,\quad
 \mu_m\mu_n^*=\mu_n^*\mu_m\quad((m,n)=1),\\
 e(r)e(t)=e(r+t),\quad e(r)^*=e(-r),\quad e(0)=1,\\
 e(r)\mu_n=\mu_ne(nr),\qquad
 \mu_ne(r)\mu_n^*=\frac1n\sum_{n\delta=r}e(\delta).
\end{gathered}
\]
The dynamics fixes $e(r)$ and sends $\mu_n$ to $n^{it}\mu_n$.
Bost and Connes prove uniqueness of KMS$_\beta$ for $0<\beta\leq1$.
Its von Neumann factor is of type $\mathrm{III}_1$; this statement concerns
the full global system, not the prime Gibbs factor.
For $\beta>1$, extreme states are indexed by
$\gamma\in\operatorname{Gal}(\mathbb Q^{\mathrm{cycl}}/\mathbb Q)$
\cite[Propositions~18, 21, 23--24, Theorem~25, and \S7]{syn-arith:BC}.
They are the normalized Gibbs states in
\[
 \pi_\gamma(\mu_m)e_n=e_{mn},\qquad
 \pi_\gamma(e(r))e_n=\gamma(e^{2\pi inr})e_n.
\]
All restrict to $\Phi_\beta$.  The additional phase observables distinguish them.

\begin{proposition}[Modular normalization]
\label{syn-arith:modular}
In the GNS representation of the prime Gibbs state, or a full
Bost--Connes Gibbs state with $\beta>1$, let $\Omega$ be the cyclic vector.
On the analytic cyclic domain the modular operator and conjugation satisfy
\[
 \Delta^{it}(a\Omega)=\alpha_{-\beta t}(a)\Omega,
 \qquad
 \Delta(\mu_m\Omega)=m^{-\beta}\mu_m\Omega,
 \qquad
 J(\mu_m\Omega)=m^{\beta/2}\mu_m^*\Omega.
\]
For a prime system, $\mu_m$ means $S_p^r$ with $m=p^r$.
\end{proposition}
\begin{proof}
For a faithful Gibbs density $D$, identify the GNS space with Hilbert--Schmidt
operators, with $\Omega=D^{1/2}$ and left multiplication as representation.
Then $\Delta(X)=DXD^{-1}$ on its maximal domain and $J(X)=X^*$.
Since $D\mu_m=m^{-\beta}\mu_mD$, these formulas give the assertions.
The von Neumann closures of the stated Gibbs representations are full operator algebras.
Moreover
\[
 \|\mu_m\Omega\|^2=1,\qquad
 \|\mu_m^*\Omega\|^2=m^{-\beta}.
\]
Thus the factor $m^{\beta/2}$ is forced by antiunitarity.
The map $a\Omega\mapsto a^*\Omega$ is the Tomita operator
$J\Delta^{1/2}$, before polar decomposition.
\end{proof}

\begin{remark}[The real Mellin transform has no heat-trace implication]
\label{syn-arith:no-gaussian-operator}
On $L^2(\mathbb R_+,dx/x)$, logarithmic coordinates identify dilation
with translation on $L^2(\mathbb R,du)$.  Its self-adjoint generator is
$P=-i\,d/du$ on $H^1(\mathbb R)$ and has spectrum $\mathbb R$.
For real $\beta>0$, $e^{-\beta P}$ is unbounded, since its Fourier multiplier
$e^{-\beta t}$ is unbounded as $t\to-\infty$.
The Gaussian Mellin identity in Proposition~\ref{syn-arith:local-tate}
therefore supplies neither this heat operator nor a trace-class gamma factor.
The meromorphic continuation of $\Lambda$ leaves
\eqref{syn-arith:trace-threshold} unchanged.
\end{remark}

\subsection{Valuation, prime support, and modified local vectors}
\label{syn-arith:two-source}

Let $\tau\geq0$ be real and set $a=p^{-\tau}$ and
\[
 \phi_{p,\tau}=\mathbf1_{\mathbb Z_p^\times}+a\mathbf1_{p\mathbb Z_p}
              =e_0+(a-1)e_1,
 \qquad e_j=\mathbf1_{p^j\mathbb Z_p}.
\]
The local Mellin and Fourier formulas are
\begin{equation}
 Z_p(\phi_{p,\tau},s)=1+\frac{ap^{-s}}{1-p^{-s}},\qquad
 \widehat\phi_{p,\tau}=e_0+(a-1)p^{-1}e_{-1}.
 \label{syn-arith:modified-tate}
\end{equation}
The first formula holds for $\operatorname{Re}s>0$.
The second follows from $\widehat e_j=p^{-j}e_{-j}$.
For $0<\operatorname{Re}s<1$ both integrals in
\[
 Z_p(\widehat\phi_{p,\tau},1-s)
 =\frac{1-p^{-s}}{1-p^{s-1}}Z_p(\phi_{p,\tau},s)
\]
converge, and the identity extends meromorphically.
Thus the local functional-equation factor is independent of $\tau$.
The vector is Fourier invariant exactly when $\tau=0$.
Its first-order defect is
$\widehat\phi_{p,\tau}-\phi_{p,\tau}
=\tau(\log p)(e_1-p^{-1}e_{-1})+O(\tau^2)$ in their finite-dimensional span.

The unitary dilation on additive $L^2(\mathbb Q_p,dx_p)$ is
$Uf(x)=p^{-1/2}f(px)$.  Indeed $dx=p\,d(px)$.
For $u_j=p^{j/2}e_j$, one has $Uu_j=u_{j-1}$,
$\mathcal Fu_j=u_{-j}$, and $\mathcal F U\mathcal F=U^{-1}$ on the even radial span of the $u_j$.
This span is invariant under these operations. On the full additive $L^2$ space, the identity is $\mathcal F U\mathcal F^{-1}=U^{-1}$, since $\mathcal F^2f(x)=f(-x)$.
It is generated by $\phi_{p,\tau}$ under $U,U^{-1},\mathcal F$:
write $\phi=u_0+bu_1$, where $b=(a-1)p^{-1/2}$.
Then $\phi-bU^{-1}\mathcal F\phi=(1-b^2)u_0$, and $|b|<1$.
For $\tau>0$, the family $(\phi_{p,\tau})_p$ changes the basic vector at every prime.
It is consequently not an elementary tensor in the algebraic restricted product
defined by the vectors $e_0$.

\begin{proposition}[Two commuting arithmetic energies]
\label{syn-arith:two-energies}
On $\ell^2(\mathbb N^\times)$ define the commuting diagonal self-adjoint operators
\[
 H_{\mathrm{val}}e_n=\log n\,e_n,\qquad
 H_{\mathrm{supp}}e_n=\log\operatorname{rad}(n)\,e_n
\]
on their respective maximal square-summability domains.
Here $\operatorname{rad}(n)=\prod_{p\mid n}p$ and $\operatorname{rad}(1)=1$.
For complex $s,\tau$, the diagonal operator
$e^{-sH_{\mathrm{val}}-\tau H_{\mathrm{supp}}}$ is trace class exactly when
\[
 \operatorname{Re}s>0,\qquad \operatorname{Re}(s+\tau)>1.
\]
In that domain its trace is
\[
 Z(s,\tau)=\sum_{n\geq1}n^{-s}\operatorname{rad}(n)^{-\tau}
 =\prod_p\left(1+\frac{p^{-s-\tau}}{1-p^{-s}}\right).
\]
For a finite nonempty set of primes, only $\operatorname{Re}s>0$ is required.
\end{proposition}
\begin{proof}
The absolute diagonal sum factors by unique factorization and Tonelli's theorem.
At any fixed prime, repeated powers require $\operatorname{Re}s>0$.
Under that condition the $p$th factor minus one is asymptotic to
$p^{-\operatorname{Re}(s+\tau)}$.  The positive Euler product is finite exactly
when the sum of these quantities over primes converges, namely when
$\operatorname{Re}(s+\tau)>1$.
The factorization of the complex series follows from absolute convergence.
\end{proof}

At one prime, put $q=p^{-s}$ for real $s>0$ and $a=p^{-\tau}>0$.
The normalized valuation law has partition function
$z_p=(1-q+aq)/(1-q)$.
For $R=v_p(n)$ and $B=\mathbf1_{R>0}$,
\[
 \mathbb E(R)=\frac{aq}{(1-q)(1-q+aq)},\qquad
 \mathbb E(B)=\frac{aq}{1-q+aq}.
\]
Differentiating $\log z_p$ multiplies these expectations by $\log p$.
At $\tau=0$ the law is geometric and
\[
 \mathbb E(R)=\frac q{1-q},\quad \mathbb E(B)=q,\quad
 \operatorname{Var}(R)=\frac q{(1-q)^2},\quad
 \operatorname{Var}(B)=q(1-q),\quad \operatorname{Cov}(R,B)=q.
\]
Hence $-\partial_s\log Z(s,0)=-\zeta'(s)/\zeta(s)$ for real $s>1$,
whereas $-\partial_\tau\log Z(s,0)=\sum_p(\log p)p^{-s}$.
The latter counts prime support once and omits higher prime powers.
The difference energy $H_{\mathrm{val}}-H_{\mathrm{supp}}$ is nonnegative.
At $\tau=0$ its expectation is
$\sum_p\sum_{r\geq2}(\log p)p^{-rs}$ for real $s>1$.
Thus the prime-power current in the explicit formula is a valuation
current, not a support current.

\begin{proposition}[Finite residue laws]
\label{syn-arith:residue-laws}
Normalize $|x|_p^{s-1}\phi_{p,\tau}(x)dx_p$ on $\mathbb Z_p$, with
real $s>0$ and $\tau\geq0$.  At $\tau=0$, its image in $\mathbb F_p$
is uniform exactly when $s=1$.  At $s=1$, write this law as $\pi$.
Then
\[
 \pi(0)=\frac a{p-1+a},\qquad
 \pi(r)=\frac1{p-1+a}\quad(r\ne0).
\]
The function $\mu(r)=1-(p-1)\pi(r)$ is another probability law.
On $L^2(\mathbb F_p)$, convolution by these laws satisfies
\[
 C_\mu+(p-1)C_\pi=pP_0,
\]
where $P_0$ projects onto constants.
On every nontrivial additive character, the respective eigenvalues are
\[
 r_p(\tau)=\frac{a-1}{p-1+a},\qquad
 d_p(\tau)=\frac{(p-1)(1-a)}{p-1+a}.
\]
Thus $C_\mu$ is positive semidefinite, while $C_\pi$ has negative eigenvalues
when $\tau>0$.
\end{proposition}
\begin{proof}
At $\tau=0$, the zero residue has mass $p^{-s}$.
Each nonzero residue has mass $(1-p^{-s})/(p-1)$.
These are equal precisely at $s=1$.
At $s=1$, each residue has additive measure $1/p$, giving the displayed law.
For $0\leq a\leq1$, all values of $\mu$ are nonnegative and their sum is one.
For a nontrivial character $\chi$, $\sum_{r\ne0}\chi(r)=-1$.
This gives both eigenvalues and proves the operator identity on a character basis.
\end{proof}

Writing the survival profile $q(r)=(p-1)\pi(r)$ as $\beta$ at zero and $B$
elsewhere gives
\[
 B=\frac{p-1}{p-1+a}=Z_p(\phi_{p,\tau},1)^{-1},\qquad
 \beta=aB,
\]
so its uniform mean is $(p-1)/p$ and its exceptional-to-background ratio is
$a$.  These two requirements determine $B,\beta$ uniquely.
The complementary law $1-q$ is the probability law $\mu$ above.

Residue and valuation information fit into an elementary fibre product.
Map $\mathbb F_p\to\{0,1\}$ by $r\mapsto\mathbf1_{r=0}$ and
$\mathbb N_0\to\{0,1\}$ by $m\mapsto\mathbf1_{m>0}$.
Then
\[
 \mathcal I_p=\mathbb F_p\times_{\{0,1\}}\mathbb N_0
 = (\mathbb F_p^\times\times\{0\})\sqcup(\{0\}\times\mathbb N_{\geq1})
\]
receives the compatible pair $(n\bmod p,v_p(n))$ for every positive integer
$n$, with the usual universal property of a fibre product of sets.
For $r\geq1$ define
\[
 \mathcal J_{p,r}=\{(a,m):a\in\mathbb Z/p^r\mathbb Z,\quad
                  m=\min(v_p(a),r)\}.
\]
Truncated valuation is well defined on the residue class, with value $r$
at the zero class.  The projection to $a$ is a bijection and the transition
maps are reduction together with $m\mapsto\min(m,r)$.
Consequently $\varprojlim_r\mathcal J_{p,r}=\mathbb Z_p$; at $a=0$ the
compatible valuations tend to $+\infty$.

For finite prime sets $S$, the algebra generated by the independent local
operators $C_{\pi_{p,\tau}}$ is
$\bigotimes_{p\in S}(\mathbb CP_{0,p}\oplus\mathbb C(1-P_{0,p}))$.
Its dimension is $2^{|S|}$.
On a product character with support $J\subseteq S$, the tensor operator has eigenvalue
\[
 \prod_{p\in J}r_p(\tau)
 =(-\tau)^{|J|}\prod_{p\in J}\frac{\log p}{p}
   +O_J(\tau^{|J|+1}).
\]

\begin{proposition}[Overlap of normalized local vectors]
\label{syn-arith:fidelity}
For real $\tau\geq0$, normalize $\phi_{p,\tau}$ in additive $L^2(\mathbb Z_p)$.
Its overlap with $\mathbf1_{\mathbb Z_p}$ is
\[
 c_p(\tau)=\frac{p-1+p^{-\tau}}
                 {\sqrt{p(p-1+p^{-2\tau})}}.
\]
For fixed $\tau>0$ there is $C_\tau>0$ such that
$\prod_{p\leq x}c_p(\tau)=C_\tau(\log x)^{-1/2}(1+o(1))$.
For fixed $t\geq0$,
\[
 \lim_{x\to\infty}\prod_{p\leq x}c_p(t/\log x)
 =\exp\left(-\frac12\int_0^1\frac{(1-e^{-tu})^2}{u}\,du\right).
\]
\end{proposition}
\begin{proof}
Integrating the vector and its square gives the formula for $c_p$.
For $a\in[0,1]$, expansion of the two logarithms gives, uniformly in $a$,
\[
 -\log\frac{p-1+a}{\sqrt{p(p-1+a^2)}}
 =\frac{(1-a)^2}{2p}+O\left(\frac{(1-a)^2}{p^2}\right).
\]
For fixed $\tau>0$, replacing $(1-p^{-\tau})^2$ by $1$ changes the prime sum
by a convergent series.  Mertens' theorem
$\sum_{p\leq x}p^{-1}=\log\log x+B+o(1)$ proves the first asymptotic
\cite{syn-arith:Mertens}.

For $a=e^{-t\log p/\log x}$, the error sum tends to zero by dominated convergence.
Partial summation of Mertens' theorem on $[x^\delta,x]$ gives
\[
 \sum_{x^\delta<p\leq x}\frac{f(\log p/\log x)}p
 \longrightarrow\int_\delta^1\frac{f(u)}u\,du.
\]
Here $f(u)=(1-e^{-tu})^2=O_t(u^2)$.
The bound $\sum_{p\leq y}(\log p)^2/p=O((\log y)^2)$ makes the omitted range
$p\leq x^\delta$ contribute $O_t(\delta^2)$.
Letting $\delta\downarrow0$ proves the asserted limit.
\end{proof}

The integral in Proposition~\ref{syn-arith:fidelity}, denoted $I(t)$, satisfies
\[
 I(t)=\frac{t^2}{2}-\frac{t^3}{3}+\frac{7t^4}{48}+O(t^5).
\]
For $t>0$, differentiation or integration of $(1-e^{-v})/v$ gives
\[
 I(t)=\gamma+\log t-\log2+2E_1(t)-E_1(2t),
 \qquad E_1(t)=\int_t^\infty e^{-v}\frac{dv}{v}.
\]
In particular the limiting overlap is $1-t^2/4+O(t^3)$ near zero
and is asymptotic to $e^{-(\gamma-\log2)/2}t^{-1/2}$ at infinity.

\subsection{Finite arithmetic configurations and sieve quotients}
\label{syn-arith:configurations}

Fix distinct integers $\mathcal H=\{h_1,\ldots,h_k\}$ with $k\geq1$.
For a positive integer $q$, define
\[
 G_q(\mathcal H)=\{a\in\mathbb Z/q\mathbb Z:
                         a+h_i\in(\mathbb Z/q\mathbb Z)^\times\text{ for all }i\}.
\]
The Chinese-remainder map restricts, for $(q,r)=1$, to
\[
 G_{qr}(\mathcal H)\xrightarrow{\sim}G_q(\mathcal H)\times G_r(\mathcal H),
 \qquad [a]\longmapsto([a]_q,[a]_r).
\]
Its two iterates for three coprime moduli agree because both record the three residues.
For primes, put $\nu_p=|\{-h_i\bmod p\}|$.
For squarefree $q$, this gives $|G_q(\mathcal H)|=\prod_{p\mid q}(p-\nu_p)$.

On $\mathscr K_S=L^2(\prod_{p\in S}\mathbb F_p)$ use uniform probability measure.
Multiplication by
$Q_{\mathcal H,p}(a)=\prod_i(1-\mathbf1_{a=-h_i})$
is an orthogonal projection.
The unit vector $1$ defines the positive functional
$\omega_S(A)=\langle1,A1\rangle$ and
\[
 \omega_S\left(\bigotimes_{p\in S}Q_{\mathcal H,p}\right)
 =\prod_{p\in S}(1-\nu_p/p).
\]
This is a factorization of finite congruence conditions, with a specified positive measure.

\begin{theorem}[Singular product and explicit tail]
\label{syn-arith:singular-series}
The set $\mathcal H$ is called admissible when $\nu_p<p$ for every prime.
For an admissible set,
\[
 \mathfrak S(\mathcal H)=\prod_p\sigma_p(\mathcal H),\qquad
 \sigma_p(\mathcal H)=\frac{1-\nu_p/p}{(1-1/p)^k}
\]
converges to a finite positive number.
Let $\Delta=\prod_{i<j}|h_i-h_j|$, with empty product $1$.
For every integer $P>2k$ exceeding all prime divisors of $\Delta$,
\[
 e^{-k^2/P}\prod_{p\leq P}\sigma_p
 \leq\mathfrak S(\mathcal H)\leq\prod_{p\leq P}\sigma_p.
\]
\end{theorem}
\begin{proof}
For $p\nmid\Delta$ the residues are distinct, so $\nu_p=k$.
For $p>k$, the convergent logarithmic expansion gives
\[
 \log\sigma_p=-\sum_{m\geq2}\frac{k^m-k}{mp^m}
             =-\frac{k(k-1)}{2p^2}+O_k(p^{-3}).
\]
Thus $\sum_p|\log\sigma_p|$ converges after the finitely many exceptional terms.
Admissibility makes those terms positive as well.
For $p>2k$, the absolute logarithm is at most
\[
 \frac12\sum_{m\geq2}(k/p)^m
 =\frac{k^2}{2p^2(1-k/p)}\leq\frac{k^2}{p^2}.
\]
Finally $\sum_{p>P}p^{-2}\leq\sum_{n>P}n^{-2}\leq P^{-1}$.
Exponentiation proves the tail inequalities.
\end{proof}

For $\mathcal H=\{0,2\}$, the factor at $2$ is $2$, and the other factors give
\[
 \mathfrak S(\{0,2\})=2\prod_{p>2}\left(1-\frac1{(p-1)^2}\right)>0.
\]
At $p=3$, the pair $\{0,6\}$ instead has $\nu_3=1$.
Thus the same one-prime Euler factor and the same number of shifts do not determine
the configuration correlation.
Positivity of either singular product is a local statement and does not prove
infinitely many prime pairs.

Choose squarefree $W$ divisible by every prime factor of $\Delta$.
Admissibility permits $v_0\in G_W(\mathcal H)$.
Let
\[
 I_N=\{n\in\mathbb Z:N\leq n<2N,\ n\equiv v_0\pmod W\}
\]
be nonempty, with $n+h_i\geq2$ for all relevant $n,i$.
Choose a finite set $\mathcal D\subset\mathbb N^k$ containing $(1,\ldots,1)$,
such that $d_1\cdots d_k$ is squarefree and coprime to $W$ for every $\boldsymbol d$.
Define the linear map and the prime-count multiplier
\[
 (T\lambda)(n)=\sum_{d_i\mid n+h_i\,\forall i}\lambda_{\boldsymbol d},\qquad
 R(n)=\sum_i\mathbf1_{\mathbb P}(n+h_i),\qquad
 (M_Rf)(n)=R(n)f(n).
\]
Here $T:\mathbb C^{\mathcal D}\to\mathbb C^{I_N}$ uses standard Euclidean inner products.
Set $A=T^*T$, $B=T^*M_RT$, and $V=(\ker T)^\perp$.

\begin{proposition}[Finite divisor-sum Rayleigh principle]
\label{syn-arith:finite-rayleigh}
The restrictions $A_V,B_V$ preserve $V$, and $A_V$ is positive definite.
Moreover
\begin{equation}
 \max_{T\lambda\ne0}\frac{\langle\lambda,B\lambda\rangle}
                           {\langle\lambda,A\lambda\rangle}
 =\lambda_{\max}(A_V^{-1/2}B_VA_V^{-1/2})
 \leq\max_{n\in I_N}R(n)\leq k.
 \label{syn-arith:rayleigh}
\end{equation}
If this quotient exceeds an integer $m\geq0$, some $n\in I_N$ has at least
$m+1$ prime translates.
For $q_i=\operatorname{lcm}(d_i,e_i)$, the entry $A_{\boldsymbol d,\boldsymbol e}$
is zero unless the $q_i$ are pairwise coprime.
In the remaining case,
\[
 \left|A_{\boldsymbol d,\boldsymbol e}-\frac{N}{W\prod_iq_i}\right|\leq1.
\]
\end{proposition}
\begin{proof}
If $p\mid q_i,q_j$ and the two congruences are satisfied, then $p\mid h_i-h_j$.
This forces $p\mid W$, contrary to the hypotheses on $\mathcal D$.
Otherwise Chinese remainders give one class modulo $W\prod_iq_i$.
Counting its intersection with $N$ consecutive integers gives the stated error.

Both $A$ and $B$ annihilate $\ker T$ and preserve its orthogonal complement.
For nonzero $v\in V$, $\langle v,Av\rangle=\|Tv\|^2>0$.
The substitution $y=A_V^{1/2}v$ proves the maximum formula by Rayleigh--Ritz.
The quotient is also
\[
 \frac{\sum_{n\in I_N}R(n)|(T\lambda)(n)|^2}
      {\sum_{n\in I_N}|(T\lambda)(n)|^2},
\]
a weighted average of integer prime counts.
This proves both bounds and the strict-threshold consequence.
\end{proof}

\begin{example}[An exact finite quotient]
Take $\mathcal H=\{0,2\}$, $W=6$, $v_0=5$, $N=20$, and
$\mathcal D=\{(1,1),(5,1),(1,5)\}$.
The interval carrier is $I_N=\{23,29,35\}$, with prime counts $(1,2,1)$.
In the indicated orders,
\[
 T=\begin{pmatrix}1&0&1\\1&0&0\\1&1&0\end{pmatrix},\quad
 A=\begin{pmatrix}3&1&1\\1&1&0\\1&0&1\end{pmatrix},\quad
 B=\begin{pmatrix}4&1&1\\1&1&0\\1&0&1\end{pmatrix}.
\]
For $\lambda=(1,-1,-1)$, $T\lambda=(0,1,0)$ and the quotient is $2$.
The generalized eigenvalues are $1,1,2$.
This calculation detects the pair $29,31$ inside the chosen finite carrier.
\end{example}

\subsection{The continuum sieve operator and its two-variable eigenfunction}
\label{syn-arith:continuum-sieve}

Let $\Delta_k=\{t_i\geq0:\sum_it_i\leq1\}$ with Lebesgue measure.
For $t_{\widehat i}\in\Delta_{k-1}$ put
$\ell_i(t_{\widehat i})=1-\sum_{j\ne i}t_j$ and define
\[
 T_i:L^2(\Delta_k)\longrightarrow L^2(\Delta_{k-1}),\qquad
 T_iF(t_{\widehat i})=\int_0^{\ell_i(t_{\widehat i})}F(t)\,dt_i.
\]
Cauchy--Schwarz gives $\|T_i\|\leq1$, and
$T_i^*g(t)=g(t_{\widehat i})$.
Hence $K_k=\sum_iT_i^*T_i$ is bounded, positive, and self-adjoint, with
\[
 M_k:=\sup_{F\ne0}\frac{\sum_i\|T_iF\|^2}{\|F\|^2}
     =\|K_k\|\leq k.
\]
For $k\geq2$ it is not compact.
Indeed, the subspace of functions constant in $t_i$ identifies with
$L^2(\Delta_{k-1},\ell_i\,dt_{\widehat i})$.
There $T_i^*T_i$ acts by multiplication by $\ell_i$ and is not compact.
Choose an orthonormal sequence supported where $\ell_i\geq1/2$.
Its $K_k$ quadratic values are at least $1/2$, contradicting compactness.

\begin{theorem}[The two-variable maximum]
\label{syn-arith:M2}
There is a unique $\lambda>1$ satisfying
\[
 \log\frac\lambda{\lambda-1}=2-\frac1\lambda.
\]
It equals $M_2$.  More explicitly, if $W_0$ is the positive real inverse of
$w\mapsto we^w$, then
\[
 M_2=\frac1{1-W_0(e^{-1})}=1.385933275998\ldots.
\]
The maximum is attained at $F(x,y)=f(x)+f(y)$, where
\[
 f(x)=\frac1{\lambda-1+x}
       +\frac1{2\lambda-1}\log\frac{\lambda-x}{\lambda-1+x},
       \qquad 0\leq x\leq1.
\]
\end{theorem}
\begin{proof}
The difference between the two sides of the scalar equation decreases
strictly from $+\infty$ to $-2$ as $\lambda$ increases from $1$ to infinity.
Its derivative is $-1/(\lambda(\lambda-1))-1/\lambda^2$.
The equation is equivalent to $we^w=e^{-1}$ with $w=1-1/\lambda$.
It also gives
\[
 \int_0^{1-x}f(y)\,dy=(\lambda-1+x)f(x).
\]
This identity follows by integrating the logarithm explicitly.
The function $f$ is positive on $[0,1/2]$ directly from its formula.
For $1/2<x<1$, the displayed integral identity makes $f(x)>0$ because $0<1-x<1/2$. The scalar equation gives $f(1)=0$. Thus $F(x,y)=f(x)+f(y)$ is strictly positive on $\Delta_2$: at least one of $x,y$ is less than one.
It follows that
\[
 K_2F(x,y)=\int_0^{1-y}F(u,y)\,du+
          \int_0^{1-x}F(x,v)\,dv=\lambda F(x,y).
\]
To prove maximality, write $h=Fu$.  Expanding the square on every coordinate fibre gives
\[
 \lambda\|h\|^2-\langle h,K_2h\rangle
 =\frac12\sum_{i=1}^2\int_{\Delta_1}
   \int_0^{\ell_i}\int_0^{\ell_i}
   F(t_{\widehat i},r)F(t_{\widehat i},s)
   |u(t_{\widehat i},r)-u(t_{\widehat i},s)|^2\,dr\,ds\,dt_{\widehat i}.
\]
The continuous positive function $F$ is bounded away from zero on the compact simplex.
The identity first holds for bounded $h$ and extends to $L^2$ by approximation.
Its right side is nonnegative, while $h=F$ gives equality.
This proves the maximum, agreeing with
\cite[Corollaries~6.2--6.3]{syn-arith:Polymath}.
\end{proof}

Maynard's sieve uses a separate distribution estimate to pass from $M_k$ to primes
\cite[Proposition~4.2]{syn-arith:Maynard}.
For $0<\vartheta<1$, suppose that for every $A_0>0$,
\[
 \sum_{q\leq x^\vartheta}\max_{(a,q)=1}
 \left|\sum_{\substack{x\leq n<2x\\n\equiv a\ (q)}}\Lambda_{\mathrm{vM}}(n)
 -\frac1{\varphi(q)}\sum_{\substack{x\leq n<2x\\(n,q)=1}}
            \Lambda_{\mathrm{vM}}(n)\right|
 \ll_{A_0,\vartheta}x(\log x)^{-A_0}.
\]
Here $\Lambda_{\mathrm{vM}}(p^r)=\log p$, and it vanishes otherwise.
If $M_k>2m/\vartheta$, every admissible $k$-tuple has infinitely many translates
containing at least $m+1$ primes
\cite[Theorem~3.8]{syn-arith:Polymath}.
The theorem about this asymptotic transfer does not follow from finite Rayleigh--Ritz alone.

The coordinate integrations are not commuting observables.  For $k=2$ write
$K^{(i)}=T_i^*T_i$, with $K^{(1)}$ integrating the first coordinate.
Direct integration gives
\[
 [K^{(1)},K^{(2)}]1=\frac{x^2-y^2}{2},\qquad
 \|[K^{(1)},K^{(2)}]1\|^2=\frac1{72},\qquad \|1\|^2=\frac12.
\]
Thus $\|[K^{(1)},K^{(2)}]\|\geq1/6$.
The norm calculation uses
$\int_{\Delta_2}x^ay^b\,dx\,dy=a!b!/(a+b+2)!$.
A finite-dimensional compression provides a lower bound, but an upper bound
requires control of its orthogonal complement.  Indeed, for a bounded
self-adjoint block operator
\[
 K=\begin{pmatrix}A&B\\B^*&C\end{pmatrix},\quad C\leq cI,\quad L>c,
\]
completion of the square proves
\[
 K\leq LI\quad\Longleftrightarrow\quad
 LI-A-B(LI-C)^{-1}B^*\geq0.
\]
Since $(LI-C)^{-1}\leq(L-c)^{-1}I$, the stronger inequality
$LI-A-BB^*/(L-c)\geq0$ suffices.  This statement separates the finite
matrix calculation from the required bound on $C$.

\subsection{Hecke characters, oscillator traces, and integral lattices}
\label{syn-arith:hecke}
Let $G=\mathrm{GL}_2(\mathbb Q_p)$, $K=\mathrm{GL}_2(\mathbb Z_p)$,
and give $G$ Haar measure with $\operatorname{vol}(K)=1$.
The spherical Hecke algebra $\mathcal H_p=C_c(K\backslash G/K,\mathbb C)$
uses convolution with this measure.  Write
$T=\mathbf1_{K\operatorname{diag}(p,1)K}$ and
$D=\mathbf1_{K\operatorname{diag}(p,p)K}$.
The normalized Satake isomorphism is
\[
 \mathcal H_p\simeq\mathbb C[z_1^{\pm1},z_2^{\pm1}]^{S_2},\qquad
 T\longmapsto p^{1/2}(z_1+z_2),\qquad D\longmapsto z_1z_2.
\]
The normalization is that of the modulus-half-power Satake transform
\cite[\S3, especially Proposition~3.6 and formula~(3.13)]{syn-arith:Gross}.
Working over $\mathbb C$ makes the chosen positive square root of $p$ explicit;
it is not an integral isomorphism with these formulas over $\mathbb Z$.

Define $H_r\in\mathcal H_p$ by $H_0=1$, $H_1=T$ and
$H_r=TH_{r-1}-pDH_{r-2}$ for $r\geq2$.
Then, as formal series,
\[
 \sum_{r\geq0}H_rX^r=(1-TX+pDX^2)^{-1}.
\]
At a Satake character $(z_1,z_2)$, put
$\alpha=p^{1/2}z_1$, $\beta=p^{1/2}z_2$.  The coefficient is
\[
 h_r(\alpha,\beta)=\sum_{j=0}^r\alpha^{r-j}\beta^j,
 \qquad \sum_{r\geq0}h_rX^r=\frac1{(1-\alpha X)(1-\beta X)}.
\]
These identities follow by multiplying two geometric series.
The recursively normalized $H_r$ need not be a single double-coset indicator:
$T*T=\mathbf1_{K\operatorname{diag}(p^2,1)K}+(p+1)D$, so that
$H_2=\mathbf1_{K\operatorname{diag}(p^2,1)K}+D$.
The coefficient $p+1$ counts index-$p$ sublattices in a rank-two lattice.

The normalized Satake parameters of a unitary representation are the $z_i$,
not the $\alpha,\beta$.  Temperedness means $|z_1|=|z_2|=1$.
For a normalized primitive holomorphic cuspidal eigenform of integer weight $w\geq2$, level $N$, and Dirichlet character $\chi$, the arithmetic normalization instead has $\alpha\beta=p^{w-1}\chi(p)$. At a prime $p\nmid N$, Deligne's bound is
$|\alpha|=|\beta|=p^{(w-1)/2}$ \cite[Theorem~8.2, p.~302]{syn-arith:Deligne}.
These normalizations cannot be interchanged with weight-zero Maass parameters.

The spherical probability measure with trivial central character can be
computed without an averaging assertion.  Let $A$ be the adjacency operator
of the $(p+1)$-regular tree and let $o$ be a vertex.
Removing its incoming edge gives a rooted branch with $p$ forward branches.
Schur complementation of $z-A$, for $\operatorname{Im}z>0$, gives the branch
and root resolvents
\[
 g(z)=\frac1{z-pg(z)}
     =\frac{z-\sqrt{z^2-4p}}{2p},\qquad
 G(z)=\langle\delta_o,(z-A)^{-1}\delta_o\rangle
     =\frac1{z-(p+1)g(z)}.
\]
The square root is the branch asymptotic to $z$ at infinity.
Taking $-\pi^{-1}\operatorname{Im}G(x+i0)$ gives the spectral density
\[
 \frac{(p+1)\sqrt{4p-x^2}}{2\pi((p+1)^2-x^2)}\,dx,
 \qquad |x|<2\sqrt p.
\]
In the variable $x=2\sqrt p\cos\theta$, it is
\[
 d\nu_p(\theta)=\frac2\pi
 \frac{p(p+1)\sin^2\theta}{p^2+1-2p\cos2\theta}\,d\theta,
 \qquad 0<\theta<\pi.
\]
The tree of homothety classes of rank-two $\mathbb Z_p$-lattices realizes
the spherical operator $T$ with trivial central character as $A$, so this
is its vacuum Plancherel measure.
It is a positive probability measure because it is a spectral measure of
the unit vector $\delta_o$, not because integration of arbitrary unitaries
against it is a projection.

\begin{proposition}[The one-mode character and the full Fock character]
\label{syn-arith:fock}
Let $\mathcal F_p$ be the rank-two bosonic Fock space with orthonormal occupation
basis indexed by nonnegative integers $n_{r,+},n_{r,-}$ of finite support,
$r\geq1$.  On their maximal diagonal domains define
\[
 E_p=(\log p)\sum_{r\geq1}r(n_{r,+}+n_{r,-}),\qquad
 Q_p=\sum_{r\geq1}(n_{r,+}-n_{r,-}).
\]
For $\theta\in\mathbb R$ and $\operatorname{Re}s>0$,
\[
 \operatorname{Tr}_{\mathcal F_p}(e^{-sE_p}e^{i\theta Q_p})
 =\prod_{r\geq1}\frac1{(1-e^{i\theta}p^{-rs})(1-e^{-i\theta}p^{-rs})}.
\]
On the subspace with $n_{r,\pm}=0$ for $r>1$, the trace is just its $r=1$ factor.
\end{proposition}
\begin{proof}
The absolute trace is bounded by
$\prod_{r\geq1}(1-p^{-r\operatorname{Re}s})^{-2}$, which converges since
$\sum_rp^{-r\operatorname{Re}s}<\infty$.  Summation over the independent
occupation numbers gives the product.  Restricting the allowed modes deletes
the corresponding factors, not their multiplicities by a regularization.
\end{proof}

For rank one, taking the restricted vacuum tensor product over primes gives
the full character $\prod_{r\geq1}\zeta(rs)$ on $\operatorname{Re}s>1$.
Its coefficient at $n=\prod_pp^{a_p}$ is
$\prod_p\mathfrak p(a_p)$, where $\mathfrak p(a)$ is the number of partitions
of $a$.  The one-mode subspace instead has coefficient one and character
$\zeta(s)$.  Hence the partition function of the prime Toeplitz system is a
one-mode character, not the character of the full Heisenberg vacuum module.
Two Heisenberg currents have Sugawara central charge two, but this fact does
not identify that central charge with any residue or pole order of $\zeta$.

There are corresponding operator obstructions.  If $L_m$ are Virasoro modes,
replacing all of them by $aL_m$ multiplies the linear term in their commutator
by $a^2$, whereas the required Virasoro relation multiplies it by $a$.
For $a\ne0,1$ this is not a conformal rescaling.
On the algebraic restricted tensor product over primes, $\sum_pL_m^{(p)}$
acts on finite-support states for $m\geq-1$, since those modes annihilate each
vacuum.  For $m\leq-2$ its value on the vacuum has infinitely many orthogonal
nonzero components of fixed positive norm.  It belongs neither to that
algebraic tensor product nor to its vacuum Hilbert completion.
An arithmetic energy weighted by $\log p$ is therefore not itself a global
Virasoro action.

Nor is a value at zero supplied by the prime zeta function
$P(s)=\sum_pp^{-s}$, initially defined for $\operatorname{Re}s>1$.
M\"obius inversion of $\log\zeta(s)=\sum_{r\geq1}P(rs)/r$ gives
\[
 P(s)=\sum_{r\geq1}\frac{\mu(r)}r\log\zeta(rs).
\]
The tail is locally normally convergent in $\operatorname{Re}s>0$ once
finitely many terms have been removed.  For each prime $q$, the term $r=q$
has a logarithmic singularity at $s=1/q$ with nonzero coefficient, and
all the other terms are locally holomorphic there after choosing logarithm
branches.  These singularities accumulate at zero, so this continuation
has no holomorphic germ there.  The value $\zeta(0)=-1/2$ belongs to a
different Dirichlet series and cannot sum a central charge indexed by primes.

\begin{example}[Adams operations and Hecke operations are different]
On $\mathbb Z[x_1,x_2,\ldots]$ define $\psi_n(x_r)=x_{nr}$ multiplicatively.
Then $\psi_m\psi_n=\psi_{mn}$, whence
$\psi_{p^2}-\psi_p^2+p\,\mathrm{id}=p\,\mathrm{id}\ne0$.
The $\mathrm{GL}_2$ recurrence with trivial diamond operator is not satisfied.
Moreover $\psi_2(x_1)-x_1^2=x_2-x_1^2$ is nonzero modulo two.
Thus these commuting endomorphisms are not Frobenius lifts on this integral
lattice.  In a torsion-free $\lambda$-ring its Adams operations must satisfy
$\psi_p(a)\equiv a^p\pmod p$ \cite{syn-arith:Wilkerson}.

This does not obstruct $\lambda$-rings from encoding Satake data.
On $\mathbb Z[z_1^{\pm1},z_2^{\pm1}]^{S_2}$ use the standard $\lambda$-ring
of virtual representations of a rank-two torus, so
$\psi_n(z_i)=z_i^n$.  Its complete symmetric functions are precisely $h_r$.
In particular, for $x=z_1+z_2$,
$\psi_2(x)=z_1^2+z_2^2$, whereas $h_2=z_1^2+z_1z_2+z_2^2$.
The distinction is between Adams and complete symmetric operations, not
between $\lambda$-rings and Hecke theory.
\end{example}

For the full rational one-prime Fock algebra $\mathbb Q[x_1,x_2,\ldots]$,
the symmetric-function lattice gives an integral form adapted to these
Adams operations.  Define $h_j$ by
\[
 \sum_{j\geq0}h_jt^j=\exp\left(\sum_{r\geq1}\frac{x_rt^r}{r}\right).
\]
Then $\mathbb Z[h_1,h_2,\ldots]$ is the ring of symmetric functions,
embedded in the rational Fock algebra, and $x_r$ represents the $r$th
power sum.  Evaluating in finitely many variables identifies
$\psi_n$ with raising each variable to its $n$th power; the congruence
$\psi_p(f)\equiv f^p\pmod p$ follows in characteristic $p$ and passes to
the stable symmetric-function ring.  Thus the integral-lattice obstruction
above is specific to $\mathbb Z[x_1,x_2,\ldots]$.  This alternative lattice
still constructs neither an arithmetic residue coupling nor a global
categorical realization.

There is a genuine one-particle Hecke action different from the Adams action.
For finite-support coefficient sequences $c(m)$, $m\geq1$, and fixed weight
$w\in\mathbb Z$, set
\[
 (\mathsf T_nc)(m)=\sum_{d\mid(n,m)}d^{w-1}c(nm/d^2).
\]
These operators preserve finite support.  On the basis $q^m$,
$\mathsf T_p=U_p+p^{w-1}V_p$, where $U_pq^m=q^{m/p}$ if $p\mid m$ and
zero otherwise, and $V_pq^m=q^{pm}$.
The identities $U_pV_p=1$ and
$V_pU_p=\mathbf1_{p\mid m}$ give
$\mathsf T_p^2=\mathsf T_{p^2}+p^{w-1}$; operators for distinct primes
commute and the divisor formula gives
$\mathsf T_m\mathsf T_n=\mathsf T_{mn}$ when $(m,n)=1$.
This defines a Hecke module, not a vertex-algebra endomorphism of the full
Fock algebra.  It mixes conformal degrees and specifies no products of
multi-particle states.

An average of unitaries is another operation that requires care.
For a probability measure $\nu$ on $[0,\pi]$ and a diagonal integer charge $Q$,
$A=\int e^{i\theta Q}\,d\nu(\theta)$ is a contraction.
On a charge-$m$ vector it acts by $\widehat\nu(m)=\int e^{im\theta}d\nu$.
It is a projection exactly when all these eigenvalues on the chosen charges
are zero or one.  A probability density positive on an interval generally
does not have this property.  In particular its charge-one eigenvalue has
positive imaginary part if it is positive on $(0,\pi)$.
Integration against a spherical Plancherel density is not automatically a
spectral projection.

\begin{example}[Finite charges and binary digits]
Retaining only exponents zero and one at each prime gives
\[
 \prod_p(1+p^{-s})=\frac{\zeta(s)}{\zeta(2s)},\qquad \operatorname{Re}s>1.
\]
Alternatively every nonnegative exponent has a unique binary expansion, giving
\[
 (1-q)^{-1}=\prod_{j\geq0}(1+q^{2^j}),\qquad |q|<1.
\]
The finite product is $(1-q^{2^{J+1}})/(1-q)$, which proves the limit.
This yields a graded vector-space identification between $\mathbb C[u]$
and a finite-support tensor product of two-dimensional spaces of weights $2^j$.
It is not an algebra identification with a tensor product of dual-number
algebras: the element in the lowest positive weight squares to zero there,
whereas $u^2\ne0$ in $\mathbb C[u]$.
\end{example}

For a finite prime set $S$, the Koszul complex on the regular sequence
$(x_p^2)_{p\in S}$ in $\mathbb C[x_p:p\in S]$ resolves the squarefree
quotient $\mathbb C[x_p:p\in S]/(x_p^2)$.
Assigning each $x_p$ bidegree $(\log p,\log p)$ gives its character
$\prod_{p\in S}(1+p^{-s-\tau})$.
At $\tau=0$, its negative logarithmic derivative has coefficients
$(-1)^{r-1}\log p$ at $p^{-rs}$; even prime-power contributions change sign.
For a smooth projective curve $X/\mathbb F_q$, the same distinction between
multiplicity and support gives the formal divisor series
\[
 \sum_{D\geq0}T^{\deg D}u^{\deg D_{\mathrm{red}}}
 =\prod_{x\in|X|}\left(1+\frac{(uT)^{\deg x}}{1-T^{\deg x}}\right).
\]
Each nonzero multiplicity at $x$ contributes one support weight and any
positive power of $T^{\deg x}$, proving the identity coefficientwise.

More generally, for a nowhere-zero scalar weight $w$ on $\mathbb N^\times$
with $w(1)=1$, the product
$e_m\star e_n=c(m,n)e_{mn}$ with $c(m,n)=w(m)w(n)/w(mn)$ is associative:
both bracketings have coefficient $w(m)w(n)w(r)/w(mnr)$.
The map $e_n\mapsto w(n)e_n$ is an algebra isomorphism to the ordinary
monoid algebra.  For $w(n)=\operatorname{rad}(n)^{-\tau}$ one has
$c(m,n)=\operatorname{rad}(\gcd(m,n))^{-\tau}$.
Such a scalar coboundary may change a chosen weighted basis or norm, but it
does not produce a new algebraic deformation class.

No finite-dimensional space with a fixed endomorphism $H$ can have
$\operatorname{Tr}(e^{-sH})=\zeta(s)$ throughout $\operatorname{Re}s>1$.
Its trace is a finite sum of exponentials and extends to an entire function;
the identity would contradict the pole of $\zeta$ at one.
This obstruction concerns an exact trace identity, not finite approximations.

\subsection{The product that a heat trace leaves undetermined}
\label{syn-arith:bar}
Use cohomological grading, so a differential has degree $+1$ and
$(V[1])^j=V^{j+1}$.  Let $A$ be an augmented dg algebra over a field $k$,
with augmentation ideal $\overline A$.  Its reduced bar coalgebra is
\[
 B(A)=\bigoplus_{n\geq0}(\overline A[1])^{\otimes n}.
\]
Write $[a]$ for the shifted element, of degree $|a|-1$.
The differential is the coderivation determined by
\[
 b_1[a]=-[da],\qquad b_2[a|b]=(-1)^{|a|}[ab].
\]
Extending these maps through a tensor word contributes the sign obtained by
moving their degree-one operator past all preceding shifted elements.
Thus the internal term on $a_i$ has sign
$-(-1)^{\sum_{j<i}(|a_j|-1)}$, and the product term on $a_i,a_{i+1}$ has sign
$(-1)^{\sum_{j<i}(|a_j|-1)+|a_i|}$.
For an ordinary augmented algebra in degree zero this reads
\[
 b[a]=0,\qquad b[a|b]=[ab],\qquad
 b[a|b|c]=[ab|c]-[a|bc].
\]
In general $b^2=0$: on length one this is $d^2=0$; on length two it is the
graded Leibniz identity; on length three it is associativity; disjoint
insertions cancel with opposite Koszul signs.  The deconcatenation coproduct
is compatible with $b$, since $b$ was defined as a coderivation.

\begin{proposition}[An Euler factor does not determine bar homology]
\label{syn-arith:two-bars}
Give $V=\bigoplus_{r\geq0}\mathbb Ce_r$ the energy $Ne_r=r\log p\,e_r$.
There are augmented algebras on $V$ with this energy as a derivation and the
same heat trace, but with inequivalent reduced bar complexes.
\end{proposition}
\begin{proof}
First take $A=\mathbb C[u]$, with $e_r=u^r$.  The free resolution
$0\to A\xrightarrow{u}A\to\mathbb C\to0$ gives
$\operatorname{Tor}^A_*(\mathbb C,\mathbb C)=\mathbb C\oplus\mathbb C\eta$
with $\eta$ in homological degree one.
Second take $A'=\mathbb C\oplus I$, where
$I=\bigoplus_{r\geq1}\mathbb Ce_r$ and $I^2=0$.
All reduced bar product terms vanish, so
$\operatorname{Tor}^{A'}_n(\mathbb C,\mathbb C)=I^{\otimes n}$ for every $n$.
Both products respect the energy grading and both heat traces are
$\sum_{r\geq0}p^{-rs}=(1-p^{-s})^{-1}$.
\end{proof}

For three polynomial prime factors the Koszul resolution gives the exterior
algebra on three degree-one classes, including a nonzero degree-three class.
There is no general concentration in degrees zero, one, and two even for this
simplest product model.  Nor does the Toeplitz algebra supply the missing
augmentation: a unital $*$-character cannot kill $S$ because $S^*S=1$.
Its characters send $S$ to a scalar of modulus one, and none is invariant
under $S\mapsto p^{it}S$.
Using a Toeplitz carrier in a bar construction therefore requires an explicit
choice of boundary modules or a different augmented algebra.

\subsection{Tame symbols and the missing coupling differential}
\label{syn-arith:gersten}
For a field $F$, $K_2^M(F)$ is the quotient of
$F^\times\otimes_{\mathbb Z}F^\times$ by $a\otimes(1-a)$ for $a\ne0,1$.
We write its group law additively and its generators as $\{a,b\}$.
For a discrete valuation $v$ with residue field $k_v$, the tame symbol is
\[
 \partial_v\{a,b\}
 =(-1)^{v(a)v(b)}\overline{a^{v(b)}b^{-v(a)}}\in k_v^\times.
\]
The expression under the bar is a unit.  It is bimultiplicative and kills
the Steinberg relation: the cases $v(a)>0$, $v(a)=0$, and $v(a)<0$ respectively
give reduction of $1-a$ to one, reduction of $a$ to one when necessary,
and cancellation of the common negative valuation and its sign.

For $F=\mathbb Q$, Tate's calculation gives the exact sequence
\[
 0\longrightarrow\mathbb Z/2\{\{-1,-1\}\}
 \longrightarrow K_2^M(\mathbb Q)
 \xrightarrow{(\partial_p)_p}\bigoplus_p\mathbb F_p^\times
 \longrightarrow0
\]
\cite[\S11]{syn-arith:MilnorK}.
In particular the two-term Gersten complex in cohomological degrees zero and
one has $H^0=\mathbb Z/2$ and $H^1=0$, and every term is torsion.
The real Hilbert symbol detects its kernel.
These facts do not assert an additive residue sum into a single common
residue field: the groups $\mathbb F_p^\times$ are different groups.
Nor is $H^1(\mathcal K_2)$ a $\mathrm{CH}^2$ formula; the usual codimension-two
Bloch formula has cohomological degree two.

\begin{example}[The diagonal symbol at an odd prime]
The Steinberg relation implies $\{a,-a\}=0$, and hence
$\{a,a\}=\{a,-1\}$, which is two-torsion.  At an odd prime,
\[
 \partial_p\{p,p\}=-1\in\mathbb F_p^\times\ne1.
\]
Consequently $\{p,p\}\ne0$ for every odd $p$.
At $p=2$, every tame symbol of $\{2,-1\}$ is trivial, and its real Hilbert
symbol is $+1$ because two is positive.  Tate's calculation therefore gives
$\{2,2\}=0$.
The relation $\{p,1-p\}=0$ cannot be substituted for $\{p,p\}=0$ in a
mixed differential calculation.
\end{example}

For contrast, reciprocity on a curve does have a common target.
For a perfect field $k$ and $f,g\in k(t)^\times$ on $\mathbb P^1_k$,
\[
 \prod_{x\in|\mathbb P^1_k|}
       N_{k(x)/k}\partial_x\{f,g\}=1.
\]
Here is an elementary proof when $k$ is algebraically closed.
By bimultiplicativity, constants and the functions $t-a$ suffice.
A constant $c$ paired with $t-a$ contributes $c$ at $a$ and $c^{-1}$ at
infinity.  For $a\ne b$, the pair $(t-a,t-b)$ contributes
$(a-b)^{-1}$ at $a$, $b-a$ at $b$, and $-1$ at infinity; their product is one.
For $a=b$, the two contributions are both $-1$.
Factoring over an algebraic closure with multiplicities and grouping
Galois orbits gives the displayed norm formula over the perfect field $k$.
This multiplicative identity is not itself a differential on a tensor
product of local oscillator spaces.

If $B$ is any complex of $\mathbb C$-vector spaces and $G$ is the two-term
torsion Gersten complex just displayed, then
$B\otimes_{\mathbb Z}G=0$ term by term.  A torsion element is killed by an
integer that is invertible on $B$.
Moreover the same vanishing holds for the derived tensor product, since
$\mathbb C$ is flat over $\mathbb Z$ and kills $G$.
Thus simply tensoring complex-linear local bars with $G$ loses all this
integral information, before any nontrivial coupling is considered.

\begin{proposition}[The equation required for a coupling]
\label{syn-arith:twisting}
Let $C$ be a coaugmented conilpotent dg coalgebra over a field $k$ and let $M$
be a cochain complex.  Give $\operatorname{End}(M)$ differential
$d(f)=d_Mf-(-1)^{|f|}fd_M$ and product composition.
For homogeneous $f,g:C\to\operatorname{End}(M)$ put
\[
 (f\star g)(c)=\sum_{(c)}(-1)^{|g||c_{(1)}|}
                     f(c_{(1)})g(c_{(2)}),
 \quad d_{\mathrm{Hom}}f=d_{\operatorname{End}}f-(-1)^{|f|}fd_C.
\]
A degree-one map $\tau$ with $\tau\eta=0$ satisfying
$d_{\mathrm{Hom}}\tau+\tau\star\tau=0$ defines a square-zero differential
on $C\otimes M$ by
\[
 D(c\otimes m)=d_Cc\otimes m+(-1)^{|c|}c\otimes d_Mm
       +\sum_{(c)}(-1)^{|c_{(1)}|}c_{(1)}\otimes\tau(c_{(2)})m.
\]
\end{proposition}
\begin{proof}
Let $D_0$ denote the first two terms and $D_\tau$ the last.
The tensor sign gives $D_0^2=0$.  The coalgebra chain identity gives
\[
 (D_0D_\tau+D_\tau D_0)(c\otimes m)
 =\sum c_{(1)}\otimes(d_{\mathrm{Hom}}\tau)(c_{(2)})m.
\]
Coassociativity and the two degree-one insertions give
$D_\tau^2(c\otimes m)=\sum c_{(1)}\otimes(\tau\star\tau)(c_{(2)})m$.
Their sum vanishes by the displayed Maurer--Cartan equation.
All sums are finite on each conilpotent element; a completed version requires
a filtration for which these maps and sums are continuous and convergent.
\end{proof}

An arithmetic realization by this mechanism would have to specify the
coalgebra on the residue data, integral local algebras or boundary modules,
the action entering $\tau$, and compatible transition maps as the set of
places grows.  An ordinary complex has no preferred coalgebra structure.
The separate identities $d_C^2=d_M^2=0$ give the uncoupled tensor
differential only; they do not give the mixed Maurer--Cartan equation.
Likewise, a diagonal energy operator with a nonzero eigenvalue cannot be a
differential, because its square has the square of that eigenvalue.

\subsection{Restricted tensor products require units and a specified base}
\label{syn-arith:categorical}
Fix a field $k$ and the symmetric monoidal category
$\mathrm{Pr}_{k,\mathrm{st}}^L$ of stable presentable $k$-linear categories,
with colimit-preserving functors and its relative tensor product over
$\operatorname{Mod}_k$.  Suppose each $\mathcal C_v$ is a presentably
symmetric monoidal category whose tensor product preserves colimits separately.
Its unit determines a colimit-preserving symmetric monoidal functor
$\operatorname{Mod}_k\to\mathcal C_v$, $V\mapsto V\otimes\mathbf1_v$.
Consequently insertion of these unit functors defines a filtered diagram of
symmetric monoidal categories indexed by finite sets of places:
\[
 \mathcal C_S=\bigotimes_{v\in S,k}\mathcal C_v,
 \qquad
 \mathcal C_S\longrightarrow\mathcal C_T,\quad
 X\longmapsto X\boxtimes\bigboxtimes_{v\in T\setminus S}\mathbf1_v.
\]
Its colimit in commutative algebra objects has the underlying filtered
colimit in $\mathrm{Pr}_{k,\mathrm{st}}^L$; tensor products preserve the
relevant colimits \cite[Proposition~3.2.3.1]{syn-arith:HA}.
This defines $\bigotimes'_v(\mathcal C_v,\mathbf1_v)$ with its compatible
symmetric monoidal structure.  The same assertion for another operad requires
that its finite tensor products and inserted unit functors form a diagram
of algebras over that operad, and that its forgetful functor creates the
filtered colimit.  This formulation includes the coherence hypotheses,
not just an identity between products of objects.

An idempotent commutative algebra $u$ with $u\otimes u\simeq u$ is not generally
the unit of $\mathcal C$.  A functor sending the source unit to $u$ cannot be
strong unital monoidal into $\mathcal C$ unless $u\simeq\mathbf1$.
One can instead use the monoidal localization of $u$-local objects, when it
exists, whose unit is $u$.  For example, $k\times0$ in
$\operatorname{Mod}_{k\times k}$ is idempotent but is not its unit $k\times k$.
This distinction is necessary before using spherical idempotents as
transition data.

For commutative $k$-algebras $R_v$, put
$R_S=\bigotimes_{v\in S,k}R_v$ and $R=\varinjlim_SR_S$.
There is an equivalence
\[
 \bigotimes'_v(\operatorname{Mod}_{R_v},R_v)
       \simeq\operatorname{Mod}_{R}.
\]
At a finite stage, the tensor product of module categories is the category
of modules over the tensor product of the algebras
\cite[Theorem~4.8.5.16]{syn-arith:HA}.
To see the filtered assertion, a colimit-preserving $k$-linear functor from
$\operatorname{Mod}_R$ to a stable presentable category $\mathcal D$ is
specified by an object of $\mathcal D$ with an $R$-action.
An $R$-action is exactly a compatible family of $R_S$-actions: all algebra
operations have finite arity.  This identifies its functor category with
the limit of the corresponding finite-stage functor categories, which is
the defining universal property of the colimit.
For the commutative spherical Hecke algebras, this constructs the global
spherical Hecke algebra and its module category.  It does not replace it
by the full, generally noncommutative, adelic Hecke algebra.

Tensor and product constructions differ even before categories enter.
The ring $\mathbb Z_2\otimes_{\mathbb Z}\mathbb Z_3$ is nonzero and is a
$\mathbb Q$-algebra: two is invertible in the second factor, three in the
first, and every other prime in both.  It identifies with
$\mathbb Q_2\otimes_{\mathbb Q}\mathbb Q_3$.
By contrast, $\mathbb Z_2\times\mathbb Z_3$ has neither two nor three invertible.
In fact the filtered tensor product over all primes inverts every prime,
whereas $\widehat{\mathbb Z}=\prod_p\mathbb Z_p$ does not.

There is a further obstruction to replacing an infinite product ring by a
product of module categories.  Let $R=\prod_{j\geq1}k$ and
$M=R/\bigoplus_{j\geq1}k$.  Then $M\ne0$, but the $j$th coordinate idempotent
$e_j$ kills $M$ for every $j$.  Thus the coordinate functor
$\operatorname{Mod}_R\to\prod_j\operatorname{Mod}_k$ kills a nonzero object.
It is not an equivalence.  The usual equivalence for finite products does
not extend to infinite products without extra conditions on modules.

An Iwahori convolution category brings a different restriction.
An $E_2$-monoidal category has a braiding, so its Grothendieck ring is
commutative.  Any proposed enhancement preserving a convolution product
whose Grothendieck ring is a noncommutative affine Hecke algebra is therefore
impossible.  The ordinary $E_1$ convolution, an $E_2$ structure on a suitable
centre, and a geometric factorization structure are distinct objects.
The category must be specified before any enhancement or global tensor
product can be claimed.

Finally, there are several notions of arithmetic configuration.
Finite subsets of the set of primes form a combinatorial indexing category;
giving that set the discrete topology is a choice, not its Zariski topology.
The closed points of $\operatorname{Spec}\mathbb Z$ inherit a cofinite Zariski
topology.  On the other hand, the relative product
$(\operatorname{Spec}\mathbb Z)^n_{/\operatorname{Spec}\mathbb Z}$ is
$\operatorname{Spec}\mathbb Z$, and its distinct-point locus is empty for
$n\geq2$.  This rules out that particular relative configuration scheme.
It proves no general nonexistence theorem for arithmetic configuration
objects over another base or in another category.
A locally constant sheaf on a totally disconnected space does not acquire
an $E_0$ factorization classification without a specified local model.

Nor does the lack of a chosen spectral coordinate on that configuration
prove a nonexistence theorem for Yang--Baxter algebras over arithmetic
rings.  On $V=\mathbb Q^2$, the flip $P$ defines the ordinary rational
matrix $R(u)=1-P/u$ over $\mathbb Q(u)$.
Expanding the two sides with the permutation relations proves
\[
 R_{12}(u-v)R_{13}(u-w)R_{23}(v-w)
 =R_{23}(v-w)R_{13}(u-w)R_{12}(u-v).
\]
This is an algebraic identity over a rational function field.
Connecting its formal parameters to finite-place collision geometry,
Frobenius operations, or a meromorphic one-form is a further construction.

For an actual small category $I$ and a covariant functor $F:I\to\mathrm{Ch}(k)$,
the simplicial replacement has $n$th term
$\bigoplus_{i_0\to\cdots\to i_n}F(i_0)$.
Its normalized total complex, with the internal differential and alternating
face maps, computes the derived colimit.  This is a well-defined incidence
bar once $I$ and $F$ are supplied.  It does not by itself construct the
collision maps of chiral homology or an arithmetic realization functor.

\subsection{The full Weil functional}
\label{syn-arith:weil}
For $a>1/2$, let $\mathcal A_a$ consist of smooth complex functions on
$\mathbb R$ for which
\[
 p_{a,j}(f)=\sup_{u\in\mathbb R}e^{a|u|}|f^{(j)}(u)|<\infty
 \quad\hbox{for every }j\geq0.
\]
Set $\mathcal A_\zeta=\bigcup_{a>1/2}\mathcal A_a$.
Use convolution, involution $f^*(u)=\overline{f(-u)}$, and Fourier transform
$\widehat f(z)=\int_{\mathbb R}f(u)e^{izu}\,du$.
Convolution and involution make $\mathcal A_\zeta$ a commutative nonunital
$*$-algebra: convolution of two functions with rate $a$ has any smaller
rate $a'>1/2$, by the triangle inequality and an integrable exponential
majorant.  Each transform is holomorphic for $|\operatorname{Im}z|<a$
and rapidly decreasing on closed smaller horizontal strips.

Write $\psi=\Gamma'/\Gamma$, and put
\[
 k_\infty(t)=\operatorname{Re}\psi\left(\frac14+\frac{it}{2}\right)-\log\pi.
\]
For $g\in\mathcal A_\zeta$, define
\begin{align}
 W(g)={}&\widehat g(i/2)+\widehat g(-i/2)
 +\frac1{2\pi}\int_{\mathbb R}\widehat g(t)k_\infty(t)\,dt \notag\\
 &-\sum_{n\geq2}\frac{\Lambda_{\mathrm{vM}}(n)}{\sqrt n}
                    \bigl(g(\log n)+g(-\log n)\bigr).
 \label{syn-arith:weil-formula}
\end{align}
All terms converge absolutely.  For example, with $a=1/2+\varepsilon$ the
prime sum is bounded by a constant times
$\sum_{n\geq2}(\log n)n^{-1-\varepsilon}$, and
$k_\infty(t)=O(\log(2+|t|))$ while $\widehat g$ decreases rapidly.

\begin{theorem}[Weil's formula and its positivity criterion]
\label{syn-arith:weil-theorem}
Let $\rho$ run over the nontrivial zeros of $\zeta$, with multiplicity, and
put $z_\rho=-i(\rho-1/2)$.  Then
\[
 W(g)=\sum_\rho\widehat g(z_\rho),\qquad g\in\mathcal A_\zeta.
\]
The series is absolutely convergent.  Moreover the Riemann hypothesis is
equivalent to $W(f^**f)\geq0$ for every $f\in\mathcal A_\zeta$.
\end{theorem}
\begin{proof}
The formula is Weil's explicit formula
\cite[formula~(11), p.~262]{syn-arith:Weil}, with the following normalization.
Put $F(s)=\widehat g(-i(s-1/2))$ and take $1<c<1+\varepsilon$.
A contour shift for $\Lambda'/\Lambda$ from $\operatorname{Re}s=c$ to
$1-c$, using $\Lambda(s)=\Lambda(1-s)$, gives
\[
 \sum_\rho F(\rho)=F(0)+F(1)
 +\frac1{2\pi i}\int_{\operatorname{Re}s=c}
       \frac{\Lambda'}{\Lambda}(s)\bigl(F(s)+F(1-s)\bigr)\,ds.
\]
The poles at zero and one have residue $-1$ in $\Lambda'/\Lambda$ and
therefore contribute the displayed positive polar terms.
Heights avoiding the zeros may be chosen so that the horizontal integrals
vanish; the rapid decrease of $F$ and the standard zero count
$N(T)=O(T\log T)$ suffice for this shift.
Insert $\zeta'/\zeta=-\sum_n\Lambda_{\mathrm{vM}}(n)n^{-s}$ on the right
line.  Fourier inversion gives the prime sum with its minus sign and
$n^{-1/2}$ weight.  The remaining logarithmic derivative is
\[
 \frac{\Gamma_{\mathbb R}'}{\Gamma_{\mathbb R}}(s)
 =-\frac12\log\pi+\frac12\psi(s/2).
\]
Moving its two terms to $\operatorname{Re}s=1/2$ gives precisely
$k_\infty$, not the expression with $+\log\pi$.
Rapid decay of $\widehat g$ on $|\operatorname{Im}z|\leq1/2$ and the zero
count prove absolute convergence of the zero sum.

For $g=f^**f$, analytic continuation of the Fourier product gives
\[
 \widehat g(z)=\widehat f(z)\overline{\widehat f(\overline z)}.
\]
Under RH every $z_\rho$ is real, so every summand is
$|\widehat f(z_\rho)|^2\geq0$.
For the converse, suppose that there is a nonreal $z_0=\gamma_0+ib$,
$b\ne0$.  Its conjugate is also a zero parameter with the same multiplicity.
The zero parameters lie in $|\operatorname{Im}z|<1/2$.
There are only finitely many distinct centered parameters
$w=z_\rho-\gamma_0$ with $|\operatorname{Re}w|\leq2$.
Choose a real polynomial $Q$ vanishing at all of them except $ib,-ib$ and
not vanishing at those two points; conjugation symmetry permits real
coefficients.  Set $P(w)=wQ(w)Q(-w)$, a real odd polynomial, and choose $f_A$
whose Fourier transform is
\[
 \widehat f_A(z)=P(z-\gamma_0)e^{-A(z-\gamma_0)^2},\qquad A>0.
\]
Its inverse transform is a polynomial times a Gaussian times a unitary
modulation, so $f_A\in\mathcal A_\zeta$.
The two selected zeros, of multiplicity $m$, contribute
\[
 -2m b^2|Q(ib)|^4e^{2Ab^2}
\]
to $W(f_A^**f_A)$, and all other zeros in the finite strip contribute zero.
On its complement, writing $w=x+iy$, one has
$x^2-y^2\geq15/4$.  For $A\geq1$, the absolute value of the remaining sum
is at most
\[
 e^{-15A/4}\sum_{|\operatorname{Re}w|>2}|P(w)|^2e^{-(x^2-1/4)}.
\]
The final sum is finite by the zero count.  The total is negative for large
$A$, contradicting positivity.  This is the separating Gaussian-polynomial
argument in Weil's lemma following formula~(11)
\cite[pp.~262--263]{syn-arith:Weil}.
\end{proof}

\begin{proposition}[The isolated archimedean term is not positive]
\label{syn-arith:gamma-negative}
The function $k_\infty$ is even, continuous, strictly increasing for $t>0$,
negative at zero, and tends to $+\infty$ as $t\to+\infty$.
Consequently it has a unique positive zero, but
$\frac1{2\pi}\int|\widehat f(t)|^2k_\infty(t)\,dt$ is negative for some
$f\in C_c^\infty(\mathbb R)$.
\end{proposition}
\begin{proof}
For $x>0$ the digamma series gives
\[
 \operatorname{Re}\psi(x+iy)=-\gamma+
 \sum_{n\geq0}\left(\frac1{n+1}-\frac{n+x}{(n+x)^2+y^2}\right).
\]
Each nonconstant summand is strictly increasing in $y>0$.
The reflection and duplication formulas for $\Gamma$ give
\[
 k_\infty(0)=-\gamma-\frac\pi2-3\log2-\log\pi<0.
\]
Stirling's formula gives
$k_\infty(t)=\log(|t|/(2\pi))+o(1)$ at infinity.
For a nonzero $\varphi\in C_c^\infty$ put
$f_R(u)=R^{-1/2}\varphi(u/R)$.  Changing variable $v=Rt$ shows that its
archimedean quadratic value tends to
\[
 \frac{k_\infty(0)}{2\pi}\int_{\mathbb R}|\widehat\varphi(v)|^2\,dv<0.
\]
Dominated convergence applies because $k_\infty(v/R)$ is bounded by a
constant times $1+\log(2+|v|)$ for $R\geq1$.
\end{proof}

For a function whose Fourier transform is supported where $k_\infty\geq0$,
the isolated integral is nonnegative.  This sufficient condition must be
interpreted in its own function space: a nonzero function in
$\mathcal A_\zeta$ cannot have its analytic Fourier transform vanish on an
open real interval.  Frequency cutoffs therefore change the test class.
They supply no positivity theorem for the complete Weil functional.

\begin{proposition}[A countable family of finite positivity tests]
\label{syn-arith:gram-tests}
Enumerate the Gaussians $g_{m,q}(u)=e^{-m(u-q)^2}$ with positive integers $m$
and rational $q$.  Form each finite Hermitian matrix
$G_{ij}=W(g_i^**g_j)$.  Then RH is equivalent to positive semidefiniteness
of every such finite matrix.
\end{proposition}
\begin{proof}
The forward implication is Theorem~\ref{syn-arith:weil-theorem}.
For the converse the span of these Gaussians is dense in the following
sense: if $f\in\mathcal A_a$ and $1/2<a'<a$, one can approximate $f$ and
any prescribed finite number of derivatives in the seminorms $p_{a',j}$.
Convolution with increasingly narrow normalized Gaussians approximates $f$
in those norms, by the strict margin $a-a'$ and uniform continuity on
compact intervals.  Truncating the convolution integral and replacing it
by Riemann sums gives finite linear combinations of Gaussian translates.
Moving the finitely many centers to rational numbers preserves the
approximation.  This proves the stated density.
The polar and prime terms of $W$ are continuous in $p_{a',0}$, and two
integrations by parts bound the archimedean term by a constant times
$p_{a',0}+p_{a',2}$.  Convolution is continuous after reducing the exponential
rate once more while keeping it greater than $1/2$.
Thus positivity on all finite Gaussian spans passes to every $f\in
\mathcal A_\zeta$, and Weil's criterion applies.
\end{proof}

Each matrix is finite, but the requirement ranges over all matrices.
There is no inference from positivity of an initial finite collection to
positivity of the full functional.  The exponential weight is also
substantive: polynomial decay in $u=\log n$ alone does not ensure convergence
of $\sum_n\Lambda_{\mathrm{vM}}(n)n^{-1/2}g(\log n)$.
For instance a smooth nonnegative function that vanishes for $u\leq1$
and equals $e^{-\sqrt u}$ for $u\geq2$ is a Schwartz function, but gives
a divergent prime sum by the prime number theorem.

For the normalized heat Gaussian $g_a(u)=(4\pi a)^{-1/2}e^{-u^2/(4a)}$,
put $h_{a,x}=g_{a/2}(\,\cdot-x)$ and $K_a(x)=W(g_a(\,\cdot-x))$.
The convolution identity $g_{a/2}*g_{a/2}=g_a$ gives the explicit kernel
matrix
\[
 W\left(\left(\sum_ic_ih_{a,x_i}\right)^**
             \left(\sum_jc_jh_{a,x_j}\right)\right)
 =\sum_{i,j}\overline c_ic_jK_a(x_j-x_i).
\]
Positivity of all these matrices for all $a>0$ and finite point sets is
also equivalent to Weil positivity, by the same Gaussian approximation.
The quantifier differs from the sieve criterion: Weil requires a
nonnegative value on every convolution square, whereas the sieve seeks
one vector with Rayleigh quotient above a threshold for an already
positive operator.  A common Hilbert-space notation does not identify
these universal and existential assertions.

\subsection{GNS construction and what an independent positive model must prove}
\label{syn-arith:gns}
Let $A$ be a possibly nonunital complex $*$-algebra and let $L:A\to\mathbb C$
be Hermitian and positive, meaning
$L(a^*)=\overline{L(a)}$ and $L(a^*a)\geq0$.
Use inner products conjugate-linear in the first variable.
The quadratic polynomial $L((a+zb)^*(a+zb))$ gives Cauchy--Schwarz,
\[
 |L(a^*b)|^2\leq L(a^*a)L(b^*b).
\]
Hence $N=\{a:L(a^*a)=0\}$ is a linear subspace, and
$\langle[a],[b]\rangle=L(a^*b)$ is a positive-definite inner product on $A/N$.

\begin{proposition}[The nonunital GNS representation]
\label{syn-arith:nonunital-gns}
The subspace $N$ is a left ideal.  Left multiplication gives a densely
defined $*$-representation $\pi$ on the Hilbert completion $\mathcal H_L$
with common invariant domain $A/N$.  Every $\pi(a)$ is closable.
It extends boundedly precisely when there is a finite $C_a$ such that
\[
 L(b^*a^*ab)\leq C_a^2L(b^*b)\quad\hbox{for all }b\in A.
\]
No cyclic vector representing $L$ is supplied by this construction in general.
\end{proposition}
\begin{proof}
If $b\in N$, Cauchy--Schwarz applied to $b$ and $a^*ab$ gives
$L(b^*a^*ab)=0$.  Thus $ab\in N$, proving the left-ideal assertion.
The identity
$\langle\pi(a)[b],[c]\rangle=\langle[b],\pi(a^*)[c]\rangle$
places the dense domain $A/N$ inside the adjoint domain of $\pi(a)$.
Therefore $\pi(a)$ is closable.  The displayed inequality is exactly the
boundedness condition on this dense domain.
If $A$ had a unit, $[1]$ would provide a cyclic vector; there is no such class
for a general nonunital algebra.
For example, let $A$ have basis $x$ with $x^*=x$ and zero multiplication,
and set $L(x)=1$.  Then $L(a^*a)=0$ for every $a$, so $\mathcal H_L=0$
although $L\ne0$.  Positivity on squares alone does not imply the additional
representability condition required for a vector state.
\end{proof}

When $A$ is a nonunital $C^*$-algebra and $L$ is a bounded positive functional,
there is a sharp unitization statement.  Set $M=\|L\|$.
The extension $\widetilde L(a+\lambda1)=L(a)+m\lambda$ on the unitization
is positive if and only if $m\geq M$.
Indeed $|L(a)|^2\leq M L(a^*a)$, proved using a positive contractive
approximate identity, makes
\[
 L(a^*a)+2\operatorname{Re}(\overline\lambda L(a))+m|\lambda|^2\geq0
\]
when $m\geq M$.  Conversely, for a positive contractive approximate identity
$e_i$, positivity at $1-e_i$ gives
$m-2L(e_i)+L(e_i^2)\geq0$; both limits of $L(e_i)$ and $L(e_i^2)$ equal
$M$, so $m\geq M$.
This argument uses a $C^*$-norm and boundedness, not just an algebraic
involution.

Applying the construction to $(\mathcal A_\zeta,W)$ is consequently
conditional on RH.  The Hermitian identity for $W$ follows directly from
\eqref{syn-arith:weil-formula}, but its positivity is exactly the assertion
of Theorem~\ref{syn-arith:weil-theorem}.  A Hilbert-space construction that
starts by assuming this positivity cannot independently establish it.

\begin{proposition}[Transfer from an independent positive realization]
\label{syn-arith:positive-transfer}
Suppose $B$ is a $*$-algebra, $\omega:B\to\mathbb C$ a Hermitian positive
functional, and $\iota:\mathcal A_\zeta\to B$ a $*$-homomorphism such that
$\omega(\iota(g))=W(g)$ for every $g\in\mathcal A_\zeta$.
Then RH holds.
The same conclusion holds for a dg $*$-algebra if these data are defined
on its degree-zero cohomology with the induced product and involution.
\end{proposition}
\begin{proof}
For each $f$,
$W(f^**f)=\omega(\iota(f)^*\iota(f))\geq0$.
Apply Theorem~\ref{syn-arith:weil-theorem}.
In the dg case the assertion is the same proof on $H^0$; compatibility of
the involution with the differential and vanishing of the functional on
boundaries must be part of the input making $H^0$ well defined.
\end{proof}

An involution is not positivity data.  On $\mathbb C^2$, for example,
the Hermitian form $|z_1|^2-|z_2|^2$ is invariant under ordinary conjugation
but has negative vectors.  An algebraic pairing or duality gives no positive
cone unless that cone and the functional's behavior on it are specified.
A chain comparison must also preserve the product and involution; an
isomorphism of graded vector spaces or equality of partition functions
does not imply the hypothesis of Proposition~\ref{syn-arith:positive-transfer}.

There is also a precise comparison of the resulting Hilbert spaces.
For a $*$-homomorphism $T:A\to B$ and $L=\omega\circ T$, the map
$[a]\mapsto[T(a)]$ preserves inner products and has kernel zero on the GNS
quotient.  It extends to an isometry $U:\mathcal H_L\to\mathcal H_\omega$
with closed range, and it intertwines left multiplication on the common
algebraic domain.  If $B$ is a $C^*$-algebra, the order inequality
\[
 \omega(d^*b^*bd)\leq\|b\|^2\omega(d^*d)
\]
gives $\|\pi_L(a)\|\leq\|T(a)\|$, so the intertwining extends to the
completed spaces.  Without that boundedness, an algebraic intertwining
identity is not an equality of arbitrary unbounded-operator domains.

Positivity need not give bounded multipliers.  For $A=\mathbb C[x]$ and
$L(p)=\pi^{-1/2}\int p(t)e^{-t^2}dt$, multiplication by $x$ is unbounded,
because its squared norm ratio on $x^n$ is $n+1/2$.
Nor does symmetry give a self-adjoint closure: on $L^2(0,\infty)$,
$-i\,d/dx$ with domain $C_c^\infty(0,\infty)$ has deficiency indices
$(1,0)$.  Solving $-if'=if$ and $-if'=-if$ gives respectively $e^{-x}$
and $e^x$, only the first square-integrable.
It therefore has no self-adjoint extension.

For an arbitrary nonunital algebra, positive extension by a unit of mass
$m$ exists exactly when $m\geq0$, $L$ is Hermitian, and
$|L(a)|^2\leq mL(a^*a)$ for every $a$; this is the same quadratic
minimization in the unit's scalar coefficient.
For $A=c_{00}(\mathbb N)$ and $L(a)=\sum_nna_n$, positivity holds but this
bound would require $n^2\leq mn$ for every point mass, which is impossible.
In the $C^*$ case the minimal mass $m=\|L\|$ has the additional property
$[e_i]\to[1]$ and gives a nondegenerate cyclic representation, since
$\|[1]-[e_i]\|^2=m-2L(e_i)+L(e_i^2)\to0$.
For larger mass the same limit is $m-\|L\|>0$
\cite[pp.~75--79]{syn-arith:Segal}.

Finally, the isotropic vectors of an indefinite Hermitian form cannot be
removed as a null subspace.  For $J=\operatorname{diag}(1,-1)$, both
$e_1+e_2$ and $e_1-e_2$ are isotropic, but their sum is not.
The radical is zero, so quotienting by it does not remove the negative
direction.  The indefinite adjoint $X^\sharp=JX^*J$ on $M_2(\mathbb C)$
and vector functional $\omega(X)=e_1^*JXe_1$ give, for $X=e_{21}$,
$\omega(X^\sharp X)=-1$.  Thus this adjoint cannot supply Weil positivity.

\subsection{Regularized determinants and the spectral implication}
\label{syn-arith:deninger}
For an operator $T$ with discrete nonzero eigenvalues $\lambda$, counted
with finite algebraic multiplicity $m_\lambda$, choose a logarithm for each
eigenvalue by a specified common spectral cut.  Suppose
\[
 \zeta_T(z)=\sum_\lambda m_\lambda e^{-z\log\lambda}
\]
converges in a right half-plane, continues meromorphically to a neighborhood
of zero, and is holomorphic there.  Define
$\det_\infty T=\exp(-\zeta_T'(0))$.
For a pencil with zero eigenvalues its determinant must be extended with
the specified zero divisor; this extension is an additional analytic
requirement, not a consequence of the definition at invertible parameters.
For positive real $c$ and compatible logarithms,
\[
 \det_\infty(cT)=c^{\zeta_T(0)}\det_\infty(T),
\]
because $\zeta_{cT}(z)=c^{-z}\zeta_T(z)$.
A scale inside a regularized determinant cannot be deleted without this
factor.

The finite-dimensional cohomological identity is unconditional.
If $C^\bullet$ is a bounded complex of finite-dimensional vector spaces and
$F$ a chain endomorphism, then, as rational functions of $T$,
\[
 \prod_j\det(1-TF\mid C^j)^{(-1)^{j+1}}
 =\prod_j\det(1-TF\mid H^j(C))^{(-1)^{j+1}}.
\]
The exact sequences through cycles and boundaries make their determinant
factors cancel.  Equivalently the logarithm is
$\sum_{r\geq1}T^r\sum_j(-1)^j\operatorname{Tr}(F^r\mid H^j(C))/r$.
This finite identity does not choose a complex from its scalar.
An acyclic summand $V\xrightarrow{1}V$ contributes one to the determinant
ratio for every endomorphism acting equally in its two degrees.
Even a single characteristic polynomial does not determine the operator:
the zero $2\times2$ matrix and a nonzero nilpotent Jordan block both have
characteristic polynomial $s^2$.

\begin{proposition}[Conditional determinant normalization]
\label{syn-arith:determinant-transfer}
Suppose complex spaces $H^0,H^1,H^2$ have endomorphisms $\Theta_i$ and
regularized determinant pencils
\[
 D_i(s)=\det_\infty\left(\frac{s-\Theta_i}{2\pi}\,\middle|\,H^i\right)
\]
that continue to single-valued meromorphic functions.
Assume $(H^0,\Theta_0)=(\mathbb C,0)$, $(H^2,\Theta_2)=(\mathbb C,1)$,
and the identity in Deninger's normalization
\[
 2^{-1/2}\Lambda(s)=\frac{D_1(s)}{D_0(s)D_2(s)}.
\]
Then
\[
 D_0(s)=\frac{s}{2\pi},\qquad D_2(s)=\frac{s-1}{2\pi},\qquad
 \xi(s)=2\sqrt2\,\pi^2D_1(s).
\]
Suppose additionally that $H^1$ has a positive Hilbert completion on which
$\Theta_1$ is densely defined, closed, and has compact resolvent, and that
its entire spectrum equals the nontrivial zero multiset of $\zeta$, with
algebraic multiplicities.  If
\[
 \mathcal D(\Theta_1^*)=\mathcal D(\Theta_1),\qquad
 \Theta_1^*=1-\Theta_1,
\]
then RH holds.
\end{proposition}
\begin{proof}
The ordinary one-dimensional determinants give $D_0,D_2$.
Multiplication by $\tfrac12s(s-1)$ proves the formula for $\xi$.
On the common adjoint domain,
$H=-i(\Theta_1-1/2)$ satisfies
$H^*=i(\Theta_1^*-1/2)=H$, with equality of domains.
Thus $H$ is self-adjoint and $\operatorname{Spec}\Theta_1\subset1/2+i\mathbb R$.
The assumed spectral equality gives RH.
\end{proof}

This is the determinant formulation of Deninger's proposed cohomological
picture \cite[\S3, equation~(3)]{syn-arith:Deninger}.
Neither a formal determinant nor a trace identity on an unspecified domain
supplies its hypotheses.  An adjoint relation on a core gives at most
skew-symmetry until essential skew-adjointness is proved.
The prime energy $He_n=(\log n)e_n$ has the logarithms of positive integers
as eigenvalues, not the zeta zero parameters; tensoring it with an infinite
degeneracy space also destroys its trace-class heat operator.
The Fock energy, this prime energy, a hypothetical self-adjoint zero-height
operator, and $\Theta_1=1/2+iH$ are four different operators.

The functional equation alone does not constrain zeros to its symmetry
line.  For $0<a<1/2$, the real polynomial
$P(s)=(s-1/2)^2-a^2$ satisfies $P(1-s)=P(s)$ but has zeros off that line.
Likewise a duality pairing need only exchange $\rho$ with $1-\overline\rho$.
Positive-definiteness and the complete adjoint-domain identity are the
additional ingredients in Proposition~\ref{syn-arith:determinant-transfer}.

Li's coefficients give another exact formulation.  With the logarithm of
$\xi$ defined near one,
\[
 \lambda_n=\left.\frac1{(n-1)!}\frac{d^n}{ds^n}
                  \bigl(s^{n-1}\log\xi(s)\bigr)\right|_{s=1}
 =\lim_{T\to\infty}\sum_{|\operatorname{Im}\rho|\leq T}
       \left[1-\left(1-\frac1\rho\right)^n\right].
\]
Li's criterion states that RH holds if and only if every $\lambda_n\geq0$
\cite{syn-arith:Li}.
Expansion of $\Gamma$ and $\zeta$ at one gives
$\lambda_1=1+\gamma/2-\tfrac12\log(4\pi)$.
To write these numbers as traces of
$1-(1-\Theta_1^{-1})^n$ requires, in addition to a spectral realization,
a regularized trace compatible with precisely the displayed symmetric
summation.  Holomorphic functional calculus alone does not define such a
trace.  Evaluating $L_0^n-(L_0-1)^n$ on a Fock space is a different
operation on a different spectrum.

\subsection{A finite-field model with an actual positive spectral form}
\label{syn-arith:finite-field}
Consider $E/\mathbb F_5$ given by $y^2=x^3-x$ and its automorphism
$I(x,y)=(-x,2y)$, with $I^2=[-1]$.
Counting the fibres above $x=0,1,2,3,4$ gives respectively
$1,1,2,2,1$ affine points and hence $\#E(\mathbb F_5)=8$.
Its Frobenius has trace $5+1-8=-2$ and characteristic polynomial
$X^2+2X+5$ on the prime-to-five Tate module.
More precisely the Frobenius endomorphism is $\pi=[-1]-[2]I$.
Indeed this endomorphism has degree five since its dual is $[-1]+[2]I$.
It acts by zero on the invariant differential $dx/y$, because
$I^*(dx/y)=2dx/y$ in characteristic five.
It is therefore a purely inseparable degree-five isogeny, hence an
automorphism followed by Frobenius.
At $P=(2,1)$, one has $I(P)=(3,2)$ and $[2]I(P)=(0,0)$, and the group law
shows $([-1]-[2]I)(P)=P$.
Among the four automorphisms $1,-1,I,-I$, only one fixes $P$, proving the
claim.

Writing $t_r=\pi^r+\overline\pi^{\,r}$, one obtains
\[
 t_0=2,\quad t_1=-2,\quad t_r=-2t_{r-1}-5t_{r-2},\qquad
 N_r:=\#E(\mathbb F_{5^r})=1+5^r-t_r.
\]
Thus $N_1,\ldots,N_6$ are
$8,32,104,640,3208,15392$.
Every closed point of degree $d\mid r$ gives $d$ rational points over
$\mathbb F_{5^r}$, so M\"obius inversion gives
$b_r=r^{-1}\sum_{d\mid r}\mu(r/d)N_d$ closed points of degree $r$.
Their first six counts are $8,12,32,152,640,2544$.
The formal exponential of the point counts is
\[
 Z_E(T)=\exp\left(\sum_{r\geq1}\frac{N_rT^r}{r}\right)
       =\frac{1+2T+5T^2}{(1-T)(1-5T)}.
\]
The Euler product over closed points counts effective divisors; comparison
of coefficients yields $\#\operatorname{Div}^{n,+}(E)(\mathbb F_5)
=2(5^n-1)$ for $n\geq1$.
These are consequences of the actual finite-field Frobenius, rather than
of prescribing a desired zero divisor.

In the basis $1,I$ of $\mathbb Z[I]$, multiplication by $\pi$ is
\[
 F=\begin{pmatrix}-1&2\\-2&-1\end{pmatrix},\qquad F^TF=5I_2.
\]
Therefore $U=F/\sqrt5$ is unitary and has eigenvalues $e^{\pm i\theta}$
with $\cos\theta=-1/\sqrt5$ and $0<\theta<\pi$.
The functional on $\mathbb C[z,z^{-1}]$ given by
$\omega(z^n)=\operatorname{Tr}(U^n)=5^{-n/2}t_n$ is positive:
$\omega(q^*q)=|q(e^{i\theta})|^2+|q(e^{-i\theta})|^2$,
where $z^*=z^{-1}$.  Its moment matrices have rank two once the polynomial
space contains two independent evaluations; its null ideal is generated
by $z^2+(2/\sqrt5)z+1$.
The mass is $\omega(1)=2$; division by two gives a state.

The local polynomial
\[
 P_E(s)=\det(1-5^{-s}F)=1+2\cdot5^{-s}+5^{1-2s}
\]
has zeros $1/2+i(\pm\theta+2\pi n)/\log5$, $n\in\mathbb Z$.
This can be realized by a genuine self-adjoint diagonal operator:
$He_n=(\theta+2\pi n)e_n$ on its maximal domain in $\ell^2(\mathbb Z)$,
$K=H\oplus(-H)$, and $\Theta=1/2+iK/\log5$.
Then $\Theta^*=1-\Theta$ on the common domain, with exactly that spectrum.
Since $H^{-2}$ is trace class, its ordinary Fredholm determinant satisfies
\[
 \det(1+a^2H^{-2})
 =\prod_{n\in\mathbb Z}\left(1+\frac{a^2}{(\theta+2\pi n)^2}\right)
 =\frac{\cosh a-\cos\theta}{1-\cos\theta}.
\]
The product converges locally uniformly in $a\in\mathbb C$; the final
identity follows from the sine product applied at $(\theta\pm ia)/2$.
This constructs a positive spectral model of one known Euler polynomial.
It supplies no operator for the complete Riemann zeta function.

\subsection{Rankin--Selberg unfolding on its convergence half-plane}
\label{syn-arith:rankin-selberg}
Let $\Gamma=\mathrm{SL}_2(\mathbb Z)$, let
$\Gamma_\infty=\{\pm\left(\begin{smallmatrix}1&n\\0&1\end{smallmatrix}\right):
n\in\mathbb Z\}$, and use $d\mu=dx\,dy/y^2$ on the upper half-plane.
For $\operatorname{Re}s>1$ define
$E(z,s)=\sum_{\gamma\in\Gamma_\infty\backslash\Gamma}
\operatorname{Im}(\gamma z)^s$.
Take a weight-zero Maass cusp form with real spectral parameter $r$ and
Fourier expansion
\[
 \phi(x+iy)=\sum_{n\ne0}\rho(n)\sqrt y\,
               K_{ir}(2\pi|n|y)e^{2\pi inx}.
\]
The hypothesis $r\in\mathbb R$ is explicit here; exceptional spectral
parameters require the corresponding convergence restrictions in the
Bessel integral.

\begin{theorem}[Unfolding and the Bessel Mellin integral]
\label{syn-arith:RS}
For $\operatorname{Re}s>1$,
\[
 \int_{\Gamma\backslash\mathbb H}|\phi(z)|^2E(z,s)\,d\mu
 =(2\pi)^{-s}G_r(s)\sum_{n\ne0}|\rho(n)|^2|n|^{-s},
\]
where, already for $\operatorname{Re}s>0$,
\[
 G_r(s)=\int_0^\infty K_{ir}(y)^2y^{s-1}dy
 =\frac{2^{s-3}\Gamma(s/2)^2\Gamma(s/2+ir)\Gamma(s/2-ir)}{\Gamma(s)}.
\]
\end{theorem}
\begin{proof}
Absolute convergence of $E$ and cusp decay justify unfolding to the strip
$0\leq x<1$, $y>0$, giving
$\int_0^\infty\int_0^1|\phi(x+iy)|^2y^{s-2}dx\,dy$.
Parseval in $x$ removes the cross terms.  Changing variable
$u=2\pi|n|y$ gives the stated Dirichlet series and $G_r$.
For real $s>1$ all terms are nonnegative, so Tonelli also justifies the
termwise integration and proves convergence; holomorphy gives the complex
half-plane statement.

For the integral use $\nu=ir$, $K_{-\nu}=K_\nu$, and
\[
 K_\nu(y)=\frac12(y/2)^\nu
       \int_0^\infty t^{-\nu-1}e^{-t-y^2/(4t)}dt.
\]
Multiplying the representations with $\nu$ and $-\nu$ and integrating $y$
first gives
\[
 2^{s-3}\Gamma(s/2)
 \int_0^\infty\int_0^\infty
 t^{s/2-\nu-1}u^{s/2+\nu-1}(t+u)^{-s/2}e^{-t-u}\,dt\,du.
\]
Set $t=av$, $u=a(1-v)$; the $a$-integral is $\Gamma(s/2)$ and the
$v$-integral is $B(s/2-\nu,s/2+\nu)$.
Absolute convergence holds for $\operatorname{Re}s>0$ because $\nu$ is
purely imaginary.  This proves the gamma expression.
\end{proof}

This is the classical Rankin--Selberg mechanism
\cite[p.~415, formula~(3)]{syn-arith:Zagier}.
For a Hecke eigenform with unitary central character, normalize the positive
Fourier coefficients by $\rho(n)=\rho(1)\lambda(n)$.
If $\alpha_p,\beta_p$ now denote the unitary-normalized parameters, then
$\lambda(p^r)=h_r(\alpha_p,\beta_p)$ and $|\alpha_p\beta_p|=1$.
At each unramified prime the formal identity
\[
 \sum_{r\geq0}|h_r(\alpha,\beta)|^2X^r
 =\frac{1-X^2}
 {(1-\alpha\overline\alpha X)(1-\alpha\overline\beta X)
  (1-\beta\overline\alpha X)(1-\beta\overline\beta X)}
\]
follows by substituting
$h_r=(\alpha^{r+1}-\beta^{r+1})/(\alpha-\beta)$ and summing four geometric
series; coincident parameters follow by a limit.
Consequently
$\sum_{n\geq1}|\lambda(n)|^2n^{-s}
=L(s,\phi\times\overline\phi)/\zeta(2s)$ in $\operatorname{Re}s>1$.
For real $s>1$ the unfolded integral is nonnegative.
At $s=1/2+it$ the continued gamma and Dirichlet factors are generally
complex, so that equality is not a positivity identity there.

Arbitrary multipliers of Fourier coefficients also do not preserve
automorphy.  For a nonzero coefficient $\rho(n_0)$ one can choose a smooth
function of $\log|n|$ supported so narrowly that only $|n|=|n_0|$ remains.
The resulting finite Fourier sum is periodic in $x$ and has the original
Laplace eigenvalue, but is not in general $\Gamma$-invariant.
An individual nonzero Bessel Fourier mode has power or logarithmic
behavior as $y\to0$, whereas invariance under $z\mapsto-1/z$ would force
the exponential cusp decay after inversion.
For a sine combination vanishing on the imaginary axis, approach zero
instead along $z=y(c+i)$ with $c\ne0$.  The finite Bessel expansion has a
nonzero power or logarithmic leading term for some such $c$, while inversion
has imaginary part $1/(y(1+c^2))$ and requires exponential decay.
Thus its norm on a strip is defined, but a Petersson norm on the quotient
and the unfolding above require a separate automorphy proof.

For Maass cusp forms over $\mathbb Q$, the bound
$|\alpha_p|,|\beta_p|\leq p^{7/64}$ is the Kim--Sarnak theorem
\cite[Appendix~2]{syn-arith:KimSarnak}; no assertion that it is optimal is
needed here.
If, additionally, the ninth symmetric power of a fixed unitary cuspidal
$\mathrm{GL}_2$ representation is cuspidal automorphic on
$\mathrm{GL}_{10}$, the Luo--Rudnick--Sarnak exponent gives
\[
 9\theta_p\leq\frac12-\frac1{101},\qquad
 \theta_p\leq\frac{11}{202}.
\]
This follows by applying the $\mathrm{GL}_{10}$ bound to the Satake
parameter $\alpha_p^9$ \cite{syn-arith:LRS}.
It is conditional on that lift for the given Maass representation.
Results about symmetric powers of holomorphic forms do not construct a
holomorphic companion with the same finite Satake parameters for an
arbitrary Maass form.

\subsection{Classical arithmetic Chern--Simons invariants and real places}
\label{syn-arith:acs}
Let $F$ be a number field, $X=\operatorname{Spec}\mathcal O_F$, and $n\geq2$.
First suppose $F$ is totally imaginary and contains all $n$th roots of
unity.  Artin--Verdier duality identifies
$H^3_{\mathrm{et}}(X,\mathbb Z/n)$ with the Pontryagin dual
$\mu_n(F)^D=\operatorname{Hom}(\mu_n(F),\mathbb Q/\mathbb Z)$.
A chosen primitive $n$th root identifies this target with
$\tfrac1n\mathbb Z/\mathbb Z$ by evaluation.
For a finite group $A$, a class $c\in H^3(A,\mathbb Z/n)$, and a continuous
$\rho:\pi_1^{\mathrm{et}}(X)\to A$, Kim's classical invariant is
\[
 \operatorname{CS}_c([\rho])=\operatorname{inv}(j^3(\rho^*c)),
\]
where $j^3$ is the natural comparison from continuous group cohomology to
\'etale cohomology \cite[\S1]{syn-arith:KimCS}.
Inner automorphisms act trivially on cohomology with these trivial
coefficients, so the invariant depends only on the conjugacy class of $\rho$.
This definition does not use a differential three-form on $X$.

At a real place, $G_v=\operatorname{Gal}(\mathbb C/\mathbb R)=C_2$ has
$H^q(G_v,\mathbb Z/2)=\mathbb Z/2$ for every $q\geq0$.
The standard periodic resolution has alternating maps $1-\sigma$ and
$1+\sigma$, both zero after applying trivial $\mathbb Z/2$ coefficients.
Thus a blanket finite two-cohomological-dimension statement cannot include
real decomposition groups.  Odd-order coefficients kill their positive
cohomology by averaging, but do not address the two-primary invariant.

There is nevertheless an actual real-place extension of the classical
construction \cite[\S\S2.2--2.3]{syn-arith:LeePark}.
Let $F_{\mathrm{un}}^f$ be the maximal extension unramified at finite places
and $F_{\mathrm{un}}$ its maximal subextension unramified also at infinity,
meaning that real places split completely.  Put
\[
 \pi=\operatorname{Gal}(F_{\mathrm{un}}^f/F),\qquad
 \pi^{\mathrm{un}}=\operatorname{Gal}(F_{\mathrm{un}}/F),\qquad
 \kappa:\pi\twoheadrightarrow\pi^{\mathrm{un}}.
\]
The compact-support cochain complex, in positive degrees relevant here, is
the mapping fibre of restriction to the real places:
\[
 C_c^r(X,M)=C^r(X,M)\oplus\prod_{v\mid\infty,\ v\text{ real}}C^{r-1}(G_v,M),
 \quad d(a,(b_v))=(da,(r_va-db_v)).
\]
Use the corresponding group-cochain mapping fibre for $\pi$.
The low-degree convention is fixed by compact-support \`etale cohomology;
the displayed ordinary cochains agree with the required local complexes in
the positive degrees used for a degree-three cocycle.
Artin--Verdier duality gives
\[
 H_c^3(X,\mathbb Z/n)\simeq\mu_n(F)^D
 \quad\text{\cite[Proposition~2.4]{syn-arith:LeePark}}.
\]
This records the target without assuming $\mu_n\subset F$.

\begin{proposition}[The compact-support lift of an everywhere unramified class]
\label{syn-arith:real-cs}
A normalized group three-cocycle $w$ on $\pi^{\mathrm{un}}$ determines
$[(\kappa^*w,(0)_v)]\in H_c^3(\pi,\mathbb Z/n)$ and hence a class in
$H_c^3(X,\mathbb Z/n)$.  For $c\in H^3(A,\mathbb Z/n)$ and
$\rho:\pi^{\mathrm{un}}\to A$, evaluation of this class defines a
conjugacy-invariant function
\[
 \operatorname{CS}_c:\operatorname{Hom}_{\mathrm{cont}}(\pi^{\mathrm{un}},A)/A
          \longrightarrow\mu_n(F)^D.
\]
\end{proposition}
\begin{proof}
Every real decomposition group maps trivially to $\pi^{\mathrm{un}}$.
A normalized cocycle vanishes when an argument is the identity; therefore
all local restrictions of $\kappa^*w$ vanish, and the pair with zero local
cochains is closed for the displayed mapping-fibre differential.
Changing $w$ by the coboundary of a normalized two-cochain changes the pair
by a mapping-fibre coboundary, for the same reason.
The group-to-\'etale comparison commutes with local restriction and therefore
induces a map of mapping fibres.
The duality invariant then gives the stated function, and conjugation
invariance follows as in Kim's construction.
This is the map $j^3_{\mathrm{un}}$ in
\cite[\S2.3, equation~(2)]{syn-arith:LeePark}.
\end{proof}

When real places exist, a primitive root of order $n>2$ cannot lie in $F$.
Thus the full $\tfrac1n\mathbb Z/\mathbb Z$ normalization with all roots
included is compatible with real places only for $n=2$.
The target $\mu_n(F)^D$ remains meaningful without that extra hypothesis.
This extension fixes a particular compact-support and real-splitting
convention; changing the infinite-place data changes the question.

For any finite groupoid $\mathcal G$ with an isomorphism-invariant action
$S:\pi_0\mathcal G\to\mathbb Q/\mathbb Z$, the scalar
\[
 Z(\mathcal G,S)=\sum_{[x]\in\pi_0\mathcal G}
             \frac{e^{2\pi iS(x)}}{|\operatorname{Aut}(x)|}
\]
is well defined and invariant under equivalence of finite groupoids.
Finiteness of the chosen arithmetic field groupoid and existence of its
action must be proved before applying this formula.
For everywhere unramified representations into a fixed finite group,
Hermite--Minkowski finiteness provides finitely many possible finite
extensions of the bounded degree in question, and the automorphism groups
are finite.  This yields finite classical sums in that setting.
It does not supply state spaces, a bordism category, or gluing laws for a
quantum theory.  Also $Z$ can vanish: two objects of trivial automorphism
group with action zero and one-half already give zero.

To define a Wilson trace at an unramified prime $\mathfrak p$, specify a
field $\sigma:\pi_1\to A$ and a complex representation $V$ of $A$.
Then $\operatorname{Tr}_V\sigma(\operatorname{Frob}_{\mathfrak p})$ is
conjugacy-invariant.  Its expectation is the weighted sum of these values
divided by $Z$, only when $Z\ne0$; evaluation at one field is not that
expectation even when the action vanishes.
For a fixed Frobenius convention and matrix $F_V$,
\[
 \det(1-TF_V)^{-1}
 =\exp\left(\sum_{r\geq1}\frac{\operatorname{Tr}(F_V^r)}rT^r\right),
 \qquad |T|\,\rho(F_V)<1.
\]
Diagonalization, or differentiation of the determinant followed by the
geometric resolvent series, proves the formula, including nonsemisimple
matrices.  It needs all power traces, not the single trace of $F_V$.
Geometric and arithmetic Frobenius are inverses; the representation and
the $L$-factor normalization must use the same choice.
A surface observable requires an additional higher holonomy construction;
an ordinary finite-group representation defines only the loop data above.

\subsection{The remaining arithmetic realization problem}
\label{syn-arith:realization-problem}
The constructions established here leave two independent tasks.
The algebraic task is to choose local augmented or boundary algebras and
their completions, construct the residue coupling of
Proposition~\ref{syn-arith:twisting}, and prove compatibility with changes
of the finite set of places.  A realization from a chiral algebra at the
archimedean place must be a specified functor into this common category;
a Gaussian Mellin transform is not such a functor.
The analytic task is to construct a positive realization of the full Weil
functional, or a determinant pencil and adjoint-domain identity with the
properties of Proposition~\ref{syn-arith:determinant-transfer}.
Neither task follows from the prime partition function.

An arithmetic Koszul-duality claim would further require an actual
augmented dg algebra $\mathcal V$, a continuous completed bar dual
$\operatorname{Hom}_{\mathrm{cont}}(B\mathcal V,k)$, and a comparison with
a specified smooth or dg Bost--Connes algebra that preserves multiplication,
differential, topology, and boundary data.  Equality of energy weights does
not give that comparison.  Iwasawa coefficients introduce derived
completion and base-change questions; vanishing of an Iwasawa $\mu$-invariant
alone gives no continuous-duality or derived-equivalence theorem.
Likewise, Hochschild homology and Hochschild cohomology are distinct
invariants, and neither their values nor a concentration range can be
read off from an Euler product.

Categorical Langlands comparisons need separately defined source and
target categories.  For example, a derived global Galois deformation
problem requires a residual representation, coefficient ring, determinant,
ramification set, and local deformation conditions.  One must construct
the resulting derived stack and its singular-support condition before
asking for an equivalence with an automorphic category.
A trace formula involving $\operatorname{Map}(B\widehat{\mathbb Z},-)$
only tests procyclic monodromy; it does not encode arbitrary inertia or
local $(\varphi,\Gamma)$-module data without extra maps and conditions.
The genuine local and global theorems in geometric settings do not
establish these mixed-characteristic, real-place, and integral comparisons
merely by a matching notation.

Finally, a profinite or continuous arithmetic Chern--Simons theory would
have to construct its limiting field spaces, measure or renormalized sum,
boundary state spaces, gluing maps, and anomaly data.  An identification
of its partition function with $\xi(s)$ must in addition construct the
parameter $s$ and prove the equality with all archimedean and polar
normalizations.  The finite classical actions of
Section~\ref{syn-arith:acs} do not contain these structures.
The scalar functional equation remains unconditional; the positive
arithmetic realization remains open.
